% paper06-semantic-moves.tex — Paper 6 of the Judgment Geometry series
\documentclass[11pt]{article}
% jugeo-common.tex — shared preamble for all Judgment Geometry papers
% Include via: % jugeo-common.tex — shared preamble for all Judgment Geometry papers
% Include via: % jugeo-common.tex — shared preamble for all Judgment Geometry papers
% Include via: \input{jugeo-common}

% ── Packages ────────────────────────────────────────────────────────────
\usepackage{amsmath,amssymb,amsthm,mathtools}
\usepackage{tikz-cd}
\usepackage{listings}
\usepackage{xcolor}
\usepackage{xspace}
\usepackage{booktabs}
\usepackage{multirow}
\usepackage{enumitem}
\usepackage[expansion=false]{microtype}
\usepackage{hyperref}
\usepackage{cleveref}
% \usepackage{thmtools}  % disabled: not available in this TeX installation
\usepackage{float}
\usepackage{caption}
\usepackage{subcaption}
\usepackage{tabularx}
\usepackage{stmaryrd}       % \llbracket, \rrbracket
\usepackage{bbm}            % \mathbbm{1}

% ── Page geometry ───────────────────────────────────────────────────────
\usepackage[margin=1in]{geometry}

% ── Theorem environments ────────────────────────────────────────────────
\theoremstyle{plain}
\newtheorem{theorem}{Theorem}[section]
\newtheorem{lemma}[theorem]{Lemma}
\newtheorem{proposition}[theorem]{Proposition}
\newtheorem{corollary}[theorem]{Corollary}

\theoremstyle{definition}
\newtheorem{definition}[theorem]{Definition}
\newtheorem{example}[theorem]{Example}
\newtheorem{construction}[theorem]{Construction}

\theoremstyle{remark}
\newtheorem{remark}[theorem]{Remark}
\newtheorem{notation}[theorem]{Notation}

% ── Python listings style ──────────────────────────────────────────────
\definecolor{jugeoBlue}{HTML}{2563EB}
\definecolor{jugeoGreen}{HTML}{059669}
\definecolor{jugeoRed}{HTML}{DC2626}
\definecolor{jugeoPurple}{HTML}{7C3AED}
\definecolor{jugeoGray}{HTML}{6B7280}
\definecolor{jugeoBg}{HTML}{F8FAFC}

\lstdefinestyle{jugeo-python}{
  language=Python,
  basicstyle=\ttfamily\small,
  keywordstyle=\color{jugeoBlue}\bfseries,
  stringstyle=\color{jugeoGreen},
  commentstyle=\color{jugeoGray}\itshape,
  numberstyle=\tiny\color{jugeoGray},
  backgroundcolor=\color{jugeoBg},
  numbers=left,
  numbersep=6pt,
  frame=single,
  framerule=0.4pt,
  rulecolor=\color{jugeoGray!40},
  breaklines=true,
  breakatwhitespace=true,
  showstringspaces=false,
  tabsize=4,
  xleftmargin=1.5em,
  framexleftmargin=1.5em,
  morekeywords={dataclass,frozen,field,Enum,IntEnum,auto,Any,
                Mapping,Sequence,Optional,match,case},
  literate={→}{{$\to$}}1 {←}{{$\leftarrow$}}1
           {≤}{{$\leq$}}1 {≥}{{$\geq$}}1 {≠}{{$\neq$}}1
           {∈}{{$\in$}}1 {∀}{{$\forall$}}1 {∃}{{$\exists$}}1
           {⊕}{{$\oplus$}}1 {⊖}{{$\ominus$}}1,
  escapeinside={(*@}{@*)},
}
\lstset{style=jugeo-python}

\lstdefinestyle{jugeo-output}{
  basicstyle=\ttfamily\small,
  backgroundcolor=\color{jugeoBg},
  frame=single,
  framerule=0.4pt,
  rulecolor=\color{jugeoGray!40},
  breaklines=true,
  showstringspaces=false,
  xleftmargin=1.5em,
  framexleftmargin=0.5em,
  numbers=none,
}

% ── JuGeo-specific macros ──────────────────────────────────────────────

% Judgment 8-tuple
\newcommand{\J}{\mathcal{J}}
\newcommand{\Jud}[8]{(#1,\,#2,\,#3,\,#4,\,#5,\,#6,\,#7,\,#8)}
\newcommand{\JudShort}[1]{\J_{#1}}

% Components
\newcommand{\coord}{\mathsf{c}}            % coordinate
\newcommand{\prop}{\varphi}                 % proposition
\newcommand{\carrier}{\mathcal{A}}          % carrier type
\newcommand{\evid}{\mathcal{E}}             % evidence bundle
\newcommand{\oblig}{\mathcal{O}}            % obligations
\newcommand{\obstr}{\mathcal{B}}            % obstructions
\newcommand{\trust}{\mathcal{T}}            % trust annotation
\newcommand{\prov}{\Pi}                     % provenance

% Trust algebra
\newcommand{\Talg}{\mathfrak{T}}            % trust ordered algebra
\newcommand{\Eadm}{\mathcal{E}_{\mathrm{adm}}}
\newcommand{\tleq}{\preceq}                 % trust ordering
\newcommand{\tjoin}{\oplus}                 % conservative join
\newcommand{\tatten}{\ominus}               % attenuation
\newcommand{\tpromote}{\uparrow_\pi}        % promotion
\newcommand{\tdemote}{\downarrow_\chi}       % demotion

% Trust tiers
\newcommand{\Tcontra}{\ensuremath{\mathsf{CONTRADICTED}}}
\newcommand{\Tunver}{\ensuremath{\mathsf{UNVERIFIED}}}
\newcommand{\Tcopilot}{\ensuremath{\mathsf{COPILOT}}}
\newcommand{\Toracle}{\ensuremath{\mathsf{ORACLE}}}
\newcommand{\Truntime}{\ensuremath{\mathsf{RUNTIME}}}
\newcommand{\Tsolver}{\ensuremath{\mathsf{SOLVER}}}
\newcommand{\Tproof}{\ensuremath{\mathsf{PROOF}}}

% Site and geometry
\newcommand{\Site}{\mathcal{S}}             % semantic site
\newcommand{\Cov}{\mathrm{Cov}}            % covering family
\newcommand{\Ob}{\mathrm{Ob}}              % objects
\newcommand{\Mor}{\mathrm{Mor}}            % morphisms
\newcommand{\res}{\mathrm{res}}            % restriction
\newcommand{\Cech}{\check{C}}              % Čech complex
\newcommand{\Hcoh}[1]{H^{#1}}             % cohomology group

% Descent
\newcommand{\descent}{\mathrm{desc}}
\newcommand{\glue}{\mathrm{glue}}
\newcommand{\overlap}[2]{#1 \cap #2}

% Semantic moves
\newcommand{\VERIFY}{\textsc{Verify}}
\newcommand{\CONSTRUCT}{\textsc{Construct}}
\newcommand{\NORMALIZE}{\textsc{Normalize}}
\newcommand{\GLUE}{\textsc{Glue}}
\newcommand{\RESTRICT}{\textsc{Restrict}}
\newcommand{\TRANSPORT}{\textsc{Transport}}
\newcommand{\REPAIR}{\textsc{Repair}}
\newcommand{\PROMOTE}{\textsc{Promote}}
\newcommand{\CHALLENGE}{\textsc{Challenge}}
\newcommand{\ESCALATE}{\textsc{Escalate}}

% Fragments
\newcommand{\QFLIA}{\texttt{QF\_LIA}}
\newcommand{\QFLRA}{\texttt{QF\_LRA}}
\newcommand{\QFBV}{\texttt{QF\_BV}}
\newcommand{\QFUF}{\texttt{QF\_UF}}

% Misc
\newcommand{\Python}{\textsc{Python}}
\newcommand{\JuGeo}{\textsc{JuGeo}}
\newcommand{\LEAN}{\textsc{Lean}\,4}
\newcommand{\Fstar}{F$^{\star}$}
\newcommand{\Coq}{\textsc{Coq}}
\newcommand{\Dafny}{\textsc{Dafny}}

% Common shortcuts
\newcommand{\ie}{i.e.\@\xspace}
\newcommand{\eg}{e.g.\@\xspace}
\newcommand{\cf}{cf.\@\xspace}
\newcommand{\wrt}{w.r.t.\@\xspace}

% Evaluation shortcuts
\newcommand{\numcases}{300}
\newcommand{\numspec}{100}
\newcommand{\numeq}{100}
\newcommand{\numbug}{100}
\newcommand{\medtime}{6.5\,ms}
\newcommand{\totaltime}{1.949\,s}


% ── Packages ────────────────────────────────────────────────────────────
\usepackage{amsmath,amssymb,amsthm,mathtools}
\usepackage{tikz-cd}
\usepackage{listings}
\usepackage{xcolor}
\usepackage{xspace}
\usepackage{booktabs}
\usepackage{multirow}
\usepackage{enumitem}
\usepackage[expansion=false]{microtype}
\usepackage{hyperref}
\usepackage{cleveref}
% \usepackage{thmtools}  % disabled: not available in this TeX installation
\usepackage{float}
\usepackage{caption}
\usepackage{subcaption}
\usepackage{tabularx}
\usepackage{stmaryrd}       % \llbracket, \rrbracket
\usepackage{bbm}            % \mathbbm{1}

% ── Page geometry ───────────────────────────────────────────────────────
\usepackage[margin=1in]{geometry}

% ── Theorem environments ────────────────────────────────────────────────
\theoremstyle{plain}
\newtheorem{theorem}{Theorem}[section]
\newtheorem{lemma}[theorem]{Lemma}
\newtheorem{proposition}[theorem]{Proposition}
\newtheorem{corollary}[theorem]{Corollary}

\theoremstyle{definition}
\newtheorem{definition}[theorem]{Definition}
\newtheorem{example}[theorem]{Example}
\newtheorem{construction}[theorem]{Construction}

\theoremstyle{remark}
\newtheorem{remark}[theorem]{Remark}
\newtheorem{notation}[theorem]{Notation}

% ── Python listings style ──────────────────────────────────────────────
\definecolor{jugeoBlue}{HTML}{2563EB}
\definecolor{jugeoGreen}{HTML}{059669}
\definecolor{jugeoRed}{HTML}{DC2626}
\definecolor{jugeoPurple}{HTML}{7C3AED}
\definecolor{jugeoGray}{HTML}{6B7280}
\definecolor{jugeoBg}{HTML}{F8FAFC}

\lstdefinestyle{jugeo-python}{
  language=Python,
  basicstyle=\ttfamily\small,
  keywordstyle=\color{jugeoBlue}\bfseries,
  stringstyle=\color{jugeoGreen},
  commentstyle=\color{jugeoGray}\itshape,
  numberstyle=\tiny\color{jugeoGray},
  backgroundcolor=\color{jugeoBg},
  numbers=left,
  numbersep=6pt,
  frame=single,
  framerule=0.4pt,
  rulecolor=\color{jugeoGray!40},
  breaklines=true,
  breakatwhitespace=true,
  showstringspaces=false,
  tabsize=4,
  xleftmargin=1.5em,
  framexleftmargin=1.5em,
  morekeywords={dataclass,frozen,field,Enum,IntEnum,auto,Any,
                Mapping,Sequence,Optional,match,case},
  literate={→}{{$\to$}}1 {←}{{$\leftarrow$}}1
           {≤}{{$\leq$}}1 {≥}{{$\geq$}}1 {≠}{{$\neq$}}1
           {∈}{{$\in$}}1 {∀}{{$\forall$}}1 {∃}{{$\exists$}}1
           {⊕}{{$\oplus$}}1 {⊖}{{$\ominus$}}1,
  escapeinside={(*@}{@*)},
}
\lstset{style=jugeo-python}

\lstdefinestyle{jugeo-output}{
  basicstyle=\ttfamily\small,
  backgroundcolor=\color{jugeoBg},
  frame=single,
  framerule=0.4pt,
  rulecolor=\color{jugeoGray!40},
  breaklines=true,
  showstringspaces=false,
  xleftmargin=1.5em,
  framexleftmargin=0.5em,
  numbers=none,
}

% ── JuGeo-specific macros ──────────────────────────────────────────────

% Judgment 8-tuple
\newcommand{\J}{\mathcal{J}}
\newcommand{\Jud}[8]{(#1,\,#2,\,#3,\,#4,\,#5,\,#6,\,#7,\,#8)}
\newcommand{\JudShort}[1]{\J_{#1}}

% Components
\newcommand{\coord}{\mathsf{c}}            % coordinate
\newcommand{\prop}{\varphi}                 % proposition
\newcommand{\carrier}{\mathcal{A}}          % carrier type
\newcommand{\evid}{\mathcal{E}}             % evidence bundle
\newcommand{\oblig}{\mathcal{O}}            % obligations
\newcommand{\obstr}{\mathcal{B}}            % obstructions
\newcommand{\trust}{\mathcal{T}}            % trust annotation
\newcommand{\prov}{\Pi}                     % provenance

% Trust algebra
\newcommand{\Talg}{\mathfrak{T}}            % trust ordered algebra
\newcommand{\Eadm}{\mathcal{E}_{\mathrm{adm}}}
\newcommand{\tleq}{\preceq}                 % trust ordering
\newcommand{\tjoin}{\oplus}                 % conservative join
\newcommand{\tatten}{\ominus}               % attenuation
\newcommand{\tpromote}{\uparrow_\pi}        % promotion
\newcommand{\tdemote}{\downarrow_\chi}       % demotion

% Trust tiers
\newcommand{\Tcontra}{\ensuremath{\mathsf{CONTRADICTED}}}
\newcommand{\Tunver}{\ensuremath{\mathsf{UNVERIFIED}}}
\newcommand{\Tcopilot}{\ensuremath{\mathsf{COPILOT}}}
\newcommand{\Toracle}{\ensuremath{\mathsf{ORACLE}}}
\newcommand{\Truntime}{\ensuremath{\mathsf{RUNTIME}}}
\newcommand{\Tsolver}{\ensuremath{\mathsf{SOLVER}}}
\newcommand{\Tproof}{\ensuremath{\mathsf{PROOF}}}

% Site and geometry
\newcommand{\Site}{\mathcal{S}}             % semantic site
\newcommand{\Cov}{\mathrm{Cov}}            % covering family
\newcommand{\Ob}{\mathrm{Ob}}              % objects
\newcommand{\Mor}{\mathrm{Mor}}            % morphisms
\newcommand{\res}{\mathrm{res}}            % restriction
\newcommand{\Cech}{\check{C}}              % Čech complex
\newcommand{\Hcoh}[1]{H^{#1}}             % cohomology group

% Descent
\newcommand{\descent}{\mathrm{desc}}
\newcommand{\glue}{\mathrm{glue}}
\newcommand{\overlap}[2]{#1 \cap #2}

% Semantic moves
\newcommand{\VERIFY}{\textsc{Verify}}
\newcommand{\CONSTRUCT}{\textsc{Construct}}
\newcommand{\NORMALIZE}{\textsc{Normalize}}
\newcommand{\GLUE}{\textsc{Glue}}
\newcommand{\RESTRICT}{\textsc{Restrict}}
\newcommand{\TRANSPORT}{\textsc{Transport}}
\newcommand{\REPAIR}{\textsc{Repair}}
\newcommand{\PROMOTE}{\textsc{Promote}}
\newcommand{\CHALLENGE}{\textsc{Challenge}}
\newcommand{\ESCALATE}{\textsc{Escalate}}

% Fragments
\newcommand{\QFLIA}{\texttt{QF\_LIA}}
\newcommand{\QFLRA}{\texttt{QF\_LRA}}
\newcommand{\QFBV}{\texttt{QF\_BV}}
\newcommand{\QFUF}{\texttt{QF\_UF}}

% Misc
\newcommand{\Python}{\textsc{Python}}
\newcommand{\JuGeo}{\textsc{JuGeo}}
\newcommand{\LEAN}{\textsc{Lean}\,4}
\newcommand{\Fstar}{F$^{\star}$}
\newcommand{\Coq}{\textsc{Coq}}
\newcommand{\Dafny}{\textsc{Dafny}}

% Common shortcuts
\newcommand{\ie}{i.e.\@\xspace}
\newcommand{\eg}{e.g.\@\xspace}
\newcommand{\cf}{cf.\@\xspace}
\newcommand{\wrt}{w.r.t.\@\xspace}

% Evaluation shortcuts
\newcommand{\numcases}{300}
\newcommand{\numspec}{100}
\newcommand{\numeq}{100}
\newcommand{\numbug}{100}
\newcommand{\medtime}{6.5\,ms}
\newcommand{\totaltime}{1.949\,s}


% ── Packages ────────────────────────────────────────────────────────────
\usepackage{amsmath,amssymb,amsthm,mathtools}
\usepackage{tikz-cd}
\usepackage{listings}
\usepackage{xcolor}
\usepackage{xspace}
\usepackage{booktabs}
\usepackage{multirow}
\usepackage{enumitem}
\usepackage[expansion=false]{microtype}
\usepackage{hyperref}
\usepackage{cleveref}
% \usepackage{thmtools}  % disabled: not available in this TeX installation
\usepackage{float}
\usepackage{caption}
\usepackage{subcaption}
\usepackage{tabularx}
\usepackage{stmaryrd}       % \llbracket, \rrbracket
\usepackage{bbm}            % \mathbbm{1}

% ── Page geometry ───────────────────────────────────────────────────────
\usepackage[margin=1in]{geometry}

% ── Theorem environments ────────────────────────────────────────────────
\theoremstyle{plain}
\newtheorem{theorem}{Theorem}[section]
\newtheorem{lemma}[theorem]{Lemma}
\newtheorem{proposition}[theorem]{Proposition}
\newtheorem{corollary}[theorem]{Corollary}

\theoremstyle{definition}
\newtheorem{definition}[theorem]{Definition}
\newtheorem{example}[theorem]{Example}
\newtheorem{construction}[theorem]{Construction}

\theoremstyle{remark}
\newtheorem{remark}[theorem]{Remark}
\newtheorem{notation}[theorem]{Notation}

% ── Python listings style ──────────────────────────────────────────────
\definecolor{jugeoBlue}{HTML}{2563EB}
\definecolor{jugeoGreen}{HTML}{059669}
\definecolor{jugeoRed}{HTML}{DC2626}
\definecolor{jugeoPurple}{HTML}{7C3AED}
\definecolor{jugeoGray}{HTML}{6B7280}
\definecolor{jugeoBg}{HTML}{F8FAFC}

\lstdefinestyle{jugeo-python}{
  language=Python,
  basicstyle=\ttfamily\small,
  keywordstyle=\color{jugeoBlue}\bfseries,
  stringstyle=\color{jugeoGreen},
  commentstyle=\color{jugeoGray}\itshape,
  numberstyle=\tiny\color{jugeoGray},
  backgroundcolor=\color{jugeoBg},
  numbers=left,
  numbersep=6pt,
  frame=single,
  framerule=0.4pt,
  rulecolor=\color{jugeoGray!40},
  breaklines=true,
  breakatwhitespace=true,
  showstringspaces=false,
  tabsize=4,
  xleftmargin=1.5em,
  framexleftmargin=1.5em,
  morekeywords={dataclass,frozen,field,Enum,IntEnum,auto,Any,
                Mapping,Sequence,Optional,match,case},
  literate={→}{{$\to$}}1 {←}{{$\leftarrow$}}1
           {≤}{{$\leq$}}1 {≥}{{$\geq$}}1 {≠}{{$\neq$}}1
           {∈}{{$\in$}}1 {∀}{{$\forall$}}1 {∃}{{$\exists$}}1
           {⊕}{{$\oplus$}}1 {⊖}{{$\ominus$}}1,
  escapeinside={(*@}{@*)},
}
\lstset{style=jugeo-python}

\lstdefinestyle{jugeo-output}{
  basicstyle=\ttfamily\small,
  backgroundcolor=\color{jugeoBg},
  frame=single,
  framerule=0.4pt,
  rulecolor=\color{jugeoGray!40},
  breaklines=true,
  showstringspaces=false,
  xleftmargin=1.5em,
  framexleftmargin=0.5em,
  numbers=none,
}

% ── JuGeo-specific macros ──────────────────────────────────────────────

% Judgment 8-tuple
\newcommand{\J}{\mathcal{J}}
\newcommand{\Jud}[8]{(#1,\,#2,\,#3,\,#4,\,#5,\,#6,\,#7,\,#8)}
\newcommand{\JudShort}[1]{\J_{#1}}

% Components
\newcommand{\coord}{\mathsf{c}}            % coordinate
\newcommand{\prop}{\varphi}                 % proposition
\newcommand{\carrier}{\mathcal{A}}          % carrier type
\newcommand{\evid}{\mathcal{E}}             % evidence bundle
\newcommand{\oblig}{\mathcal{O}}            % obligations
\newcommand{\obstr}{\mathcal{B}}            % obstructions
\newcommand{\trust}{\mathcal{T}}            % trust annotation
\newcommand{\prov}{\Pi}                     % provenance

% Trust algebra
\newcommand{\Talg}{\mathfrak{T}}            % trust ordered algebra
\newcommand{\Eadm}{\mathcal{E}_{\mathrm{adm}}}
\newcommand{\tleq}{\preceq}                 % trust ordering
\newcommand{\tjoin}{\oplus}                 % conservative join
\newcommand{\tatten}{\ominus}               % attenuation
\newcommand{\tpromote}{\uparrow_\pi}        % promotion
\newcommand{\tdemote}{\downarrow_\chi}       % demotion

% Trust tiers
\newcommand{\Tcontra}{\ensuremath{\mathsf{CONTRADICTED}}}
\newcommand{\Tunver}{\ensuremath{\mathsf{UNVERIFIED}}}
\newcommand{\Tcopilot}{\ensuremath{\mathsf{COPILOT}}}
\newcommand{\Toracle}{\ensuremath{\mathsf{ORACLE}}}
\newcommand{\Truntime}{\ensuremath{\mathsf{RUNTIME}}}
\newcommand{\Tsolver}{\ensuremath{\mathsf{SOLVER}}}
\newcommand{\Tproof}{\ensuremath{\mathsf{PROOF}}}

% Site and geometry
\newcommand{\Site}{\mathcal{S}}             % semantic site
\newcommand{\Cov}{\mathrm{Cov}}            % covering family
\newcommand{\Ob}{\mathrm{Ob}}              % objects
\newcommand{\Mor}{\mathrm{Mor}}            % morphisms
\newcommand{\res}{\mathrm{res}}            % restriction
\newcommand{\Cech}{\check{C}}              % Čech complex
\newcommand{\Hcoh}[1]{H^{#1}}             % cohomology group

% Descent
\newcommand{\descent}{\mathrm{desc}}
\newcommand{\glue}{\mathrm{glue}}
\newcommand{\overlap}[2]{#1 \cap #2}

% Semantic moves
\newcommand{\VERIFY}{\textsc{Verify}}
\newcommand{\CONSTRUCT}{\textsc{Construct}}
\newcommand{\NORMALIZE}{\textsc{Normalize}}
\newcommand{\GLUE}{\textsc{Glue}}
\newcommand{\RESTRICT}{\textsc{Restrict}}
\newcommand{\TRANSPORT}{\textsc{Transport}}
\newcommand{\REPAIR}{\textsc{Repair}}
\newcommand{\PROMOTE}{\textsc{Promote}}
\newcommand{\CHALLENGE}{\textsc{Challenge}}
\newcommand{\ESCALATE}{\textsc{Escalate}}

% Fragments
\newcommand{\QFLIA}{\texttt{QF\_LIA}}
\newcommand{\QFLRA}{\texttt{QF\_LRA}}
\newcommand{\QFBV}{\texttt{QF\_BV}}
\newcommand{\QFUF}{\texttt{QF\_UF}}

% Misc
\newcommand{\Python}{\textsc{Python}}
\newcommand{\JuGeo}{\textsc{JuGeo}}
\newcommand{\LEAN}{\textsc{Lean}\,4}
\newcommand{\Fstar}{F$^{\star}$}
\newcommand{\Coq}{\textsc{Coq}}
\newcommand{\Dafny}{\textsc{Dafny}}

% Common shortcuts
\newcommand{\ie}{i.e.\@\xspace}
\newcommand{\eg}{e.g.\@\xspace}
\newcommand{\cf}{cf.\@\xspace}
\newcommand{\wrt}{w.r.t.\@\xspace}

% Evaluation shortcuts
\newcommand{\numcases}{300}
\newcommand{\numspec}{100}
\newcommand{\numeq}{100}
\newcommand{\numbug}{100}
\newcommand{\medtime}{6.5\,ms}
\newcommand{\totaltime}{1.949\,s}

% comprehensive-data.tex — AUTO-GENERATED from real jugeo experiments
% DO NOT EDIT — regenerate with: python3 experiments/run_comprehensive_experiments.py
% Generated: 2026-03-23 09:19:06

% --- Size scaling metrics ---
\newcommand{\compTotalPrograms}{18}
\newcommand{\compTotalVerified}{18}
\newcommand{\compOverallAccuracy}{100.0\%}
\newcommand{\compTinyCount}{5}
\newcommand{\compTinyVerified}{5}
\newcommand{\compTinyMeanCoords}{3}
\newcommand{\compTinyMeanProps}{6}
\newcommand{\compTinyMeanTime}{4.95\,s}
\newcommand{\compTinyMeanLines}{5}
\newcommand{\compTinyMinTime}{4.45\,s}
\newcommand{\compTinyMaxTime}{5.51\,s}
\newcommand{\compSmallCount}{5}
\newcommand{\compSmallVerified}{5}
\newcommand{\compSmallMeanCoords}{3.6}
\newcommand{\compSmallMeanProps}{7.2}
\newcommand{\compSmallMeanTime}{4.85\,s}
\newcommand{\compSmallMeanLines}{16.0}
\newcommand{\compSmallMinTime}{4.30\,s}
\newcommand{\compSmallMaxTime}{5.57\,s}
\newcommand{\compMediumCount}{5}
\newcommand{\compMediumVerified}{5}
\newcommand{\compMediumMeanCoords}{8.6}
\newcommand{\compMediumMeanProps}{17.2}
\newcommand{\compMediumMeanTime}{4.47\,s}
\newcommand{\compMediumMeanLines}{45.0}
\newcommand{\compMediumMinTime}{4.26\,s}
\newcommand{\compMediumMaxTime}{4.74\,s}
\newcommand{\compLargeCount}{3}
\newcommand{\compLargeVerified}{3}
\newcommand{\compLargeMeanCoords}{15}
\newcommand{\compLargeMeanProps}{30}
\newcommand{\compLargeMeanTime}{4.31\,s}
\newcommand{\compLargeMeanLines}{79.0}
\newcommand{\compLargeMinTime}{4.27\,s}
\newcommand{\compLargeMaxTime}{4.38\,s}
% --- Aggregate timing ---
\newcommand{\compTimeMean}{4.68\,s}
\newcommand{\compTimeMedian}{4.50\,s}
\newcommand{\compTimeMin}{4.265\,s}
\newcommand{\compTimeMax}{5.571\,s}
\newcommand{\compTimeTotal}{84.3\,s}
\newcommand{\compTimeStdev}{0.45\,s}
% --- Aggregate structure ---
\newcommand{\compCoordsMean}{6.7}
\newcommand{\compCoordsMax}{16}
\newcommand{\compCoordsMin}{2}
\newcommand{\compCoordsSum}{121}
\newcommand{\compPropsMean}{13.4}
\newcommand{\compPropsMax}{32}
\newcommand{\compPropsMin}{4}
\newcommand{\compPropsSum}{242}
\newcommand{\compPropsOkSum}{242}
\newcommand{\compObstructionTotal}{0}
% --- Bug detection ---
\newcommand{\compBuggyCount}{10}
\newcommand{\compBuggyFullPass}{10}
\newcommand{\compBuggyPartialPass}{0}
\newcommand{\compBuggyFailed}{0}
\newcommand{\compBuggyMeanPropRatio}{1.00}
\newcommand{\compCorrectMeanPropRatio}{1.00}
\newcommand{\compBuggyMeanTime}{5.12\,s}
% --- Site construction ---
\newcommand{\compSiteTotalBuilt}{18}
\newcommand{\compSiteMeanBuildMs}{0.05}
\newcommand{\compSiteMaxBuildMs}{0.14}
\newcommand{\compSiteMeanCoords}{6.5}
\newcommand{\compSiteMeanMorphisms}{14.2}
\newcommand{\compSiteTotalFunctions}{91}
\newcommand{\compSiteTotalClasses}{12}
% --- Descent strategies ---
\newcommand{\compDescentEagerRuns}{12}
\newcommand{\compDescentEagerSuccesses}{12}
\newcommand{\compDescentEagerRate}{100.0\%}
\newcommand{\compDescentEagerMeanMs}{0.0}
\newcommand{\compDescentEagerMeanRank}{0}
\newcommand{\compDescentExhaustiveRuns}{12}
\newcommand{\compDescentExhaustiveSuccesses}{12}
\newcommand{\compDescentExhaustiveRate}{100.0\%}
\newcommand{\compDescentExhaustiveMeanMs}{0.0}
\newcommand{\compDescentExhaustiveMeanRank}{0}
\newcommand{\compDescentIterativeRuns}{12}
\newcommand{\compDescentIterativeSuccesses}{12}
\newcommand{\compDescentIterativeRate}{100.0\%}
\newcommand{\compDescentIterativeMeanMs}{0.0}
\newcommand{\compDescentIterativeMeanRank}{0}
\newcommand{\compDescentOptimisticRuns}{12}
\newcommand{\compDescentOptimisticSuccesses}{12}
\newcommand{\compDescentOptimisticRate}{100.0\%}
\newcommand{\compDescentOptimisticMeanMs}{0.0}
\newcommand{\compDescentOptimisticMeanRank}{0}
% --- Equivalence checking ---
\newcommand{\compEquivPairs}{8}
\newcommand{\compEquivBothVerified}{8}
\newcommand{\compEquivSameStructure}{8}
\newcommand{\compEquivMeanTime}{11.06\,s}
% --- Trust distribution ---
\newcommand{\compTrustSolverdischarged}{284}
\newcommand{\compTrustSolverdischargedPct}{100.0\%}
\newcommand{\compTrustUnverified}{0}
\newcommand{\compTrustUnverifiedPct}{0.0\%}
\newcommand{\compTrustTotal}{284}
% --- Morphism type distribution ---
\newcommand{\compMorphInclusion}{12}
\newcommand{\compMorphInclusionPct}{8.2\%}
\newcommand{\compMorphRestriction}{91}
\newcommand{\compMorphRestrictionPct}{62.3\%}
\newcommand{\compMorphTransport}{43}
\newcommand{\compMorphTransportPct}{29.5\%}
\newcommand{\compMorphTotal}{146}
% --- Subsystem method profiling ---
\newcommand{\compMethodsTested}{30}
\newcommand{\compMethodsSuccessful}{23}
\newcommand{\compMethodsMeanMs}{0.02}
\newcommand{\compProfAnalogyTransportMeanMs}{0.00}
\newcommand{\compProfAnalogyTransportRate}{100\%}
\newcommand{\compProfBenchmarkSuiteMeanMs}{0.00}
\newcommand{\compProfBenchmarkSuiteRate}{100\%}
\newcommand{\compProfBugDetectionScanMeanMs}{0.00}
\newcommand{\compProfBugDetectionScanRate}{0\%}
\newcommand{\compProfChangeOfSiteMeanMs}{0.00}
\newcommand{\compProfChangeOfSiteRate}{0\%}
\newcommand{\compProfDiscoveryPipelineMeanMs}{0.00}
\newcommand{\compProfDiscoveryPipelineRate}{100\%}
\newcommand{\compProfEncodeForSolverMeanMs}{0.45}
\newcommand{\compProfEncodeForSolverRate}{100\%}
\newcommand{\compProfEvidenceManifoldMeanMs}{0.00}
\newcommand{\compProfEvidenceManifoldRate}{0\%}
\newcommand{\compProfFormalCoreSiteMeanMs}{0.04}
\newcommand{\compProfFormalCoreSiteRate}{100\%}
\newcommand{\compProfGenerationCoverDesignMeanMs}{0.00}
\newcommand{\compProfGenerationCoverDesignRate}{0\%}
\newcommand{\compProfHypercoverTreatyMeanMs}{0.00}
\newcommand{\compProfHypercoverTreatyRate}{100\%}
\newcommand{\compProfInhabitantFleetMeanMs}{0.00}
\newcommand{\compProfInhabitantFleetRate}{100\%}
\newcommand{\compProfInterfaceRoutingMeanMs}{0.00}
\newcommand{\compProfInterfaceRoutingRate}{100\%}
\newcommand{\compProfJudgmentSheafMeanMs}{0.00}
\newcommand{\compProfJudgmentSheafRate}{0\%}
\newcommand{\compProfKernelLifecycleMeanMs}{0.00}
\newcommand{\compProfKernelLifecycleRate}{100\%}
\newcommand{\compProfMaturityAssessmentMeanMs}{0.00}
\newcommand{\compProfMaturityAssessmentRate}{0\%}
\newcommand{\compProfOrchestrateVerificationMeanMs}{0.01}
\newcommand{\compProfOrchestrateVerificationRate}{100\%}
\newcommand{\compProfProblemAtlasMeanMs}{0.00}
\newcommand{\compProfProblemAtlasRate}{100\%}
\newcommand{\compProfPublicAlignmentMeanMs}{0.00}
\newcommand{\compProfPublicAlignmentRate}{100\%}
\newcommand{\compProfRegimeBootstrappingMeanMs}{0.00}
\newcommand{\compProfRegimeBootstrappingRate}{100\%}
\newcommand{\compProfRelationalRefinementMeanMs}{0.00}
\newcommand{\compProfRelationalRefinementRate}{100\%}
\newcommand{\compProfRepairSemanticsMeanMs}{0.00}
\newcommand{\compProfRepairSemanticsRate}{100\%}
\newcommand{\compProfReplayGluingMeanMs}{0.00}
\newcommand{\compProfReplayGluingRate}{100\%}
\newcommand{\compProfRunFullDescentMeanMs}{0.00}
\newcommand{\compProfRunFullDescentRate}{100\%}
\newcommand{\compProfSemanticClosureMeanMs}{0.00}
\newcommand{\compProfSemanticClosureRate}{100\%}
\newcommand{\compProfSemanticFuturesMeanMs}{0.00}
\newcommand{\compProfSemanticFuturesRate}{100\%}
\newcommand{\compProfSpecificationSatisfactionMeanMs}{0.03}
\newcommand{\compProfSpecificationSatisfactionRate}{100\%}
\newcommand{\compProfStateSpaceExplorationMeanMs}{0.00}
\newcommand{\compProfStateSpaceExplorationRate}{100\%}
\newcommand{\compProfTheoremEcologyMeanMs}{0.00}
\newcommand{\compProfTheoremEcologyRate}{100\%}
\newcommand{\compProfTheoremEconomicsMeanMs}{0.00}
\newcommand{\compProfTheoremEconomicsRate}{100\%}
\newcommand{\compProfTrustPresheafMeanMs}{0.00}
\newcommand{\compProfTrustPresheafRate}{0\%}


% ── Additional packages for this paper ────────────────────────────────
\usepackage{mathrsfs}
\usepackage{array}
\microtypesetup{expansion=false}

% ── Paper-local macros ────────────────────────────────────────────────
\newcommand{\SM}{\mathsf{SM}}
\newcommand{\PS}{\mathsf{PS}}
\newcommand{\PScat}{\mathbf{PS}}
\newcommand{\Mcat}{\mathbf{M}}
\newcommand{\pre}{\mathrm{pre}}
\newcommand{\post}{\mathrm{post}}
\newcommand{\gain}{\mathrm{gain}}
\newcommand{\cost}{\mathrm{cost}}
\newcommand{\floor}{\mathrm{floor}}
\newcommand{\kind}{\mathrm{kind}}
\newcommand{\target}{\mathrm{tgt}}
\newcommand{\norm}[1]{\lVert #1 \rVert}
\newcommand{\obstrnorm}{\norm{\obstr}}
\newcommand{\MoveKind}{\mathsf{MoveKind}}
\newcommand{\RuleKind}{\mathsf{RuleKind}}
\newcommand{\TransKind}{\mathsf{TransKind}}
\newcommand{\CtrlStrat}{\mathsf{CtrlStrat}}
\newcommand{\ConstrStatus}{\mathsf{ConstrStatus}}
\newcommand{\SrcChan}{\mathsf{SrcChan}}
\newcommand{\deltaS}{\Delta S}
\newcommand{\deltaO}{\Delta O}
\newcommand{\deltaE}{\Delta E}
\newcommand{\deltaX}{\Delta X}
\newcommand{\deltaK}{\Delta K}

\title{\textbf{Proof Search Beyond Term Rewriting}}

\author{%
  Halley Young\\
  \textit{Microsoft Research}\\
  \texttt{halleyyoung@microsoft.com}
}
\date{}

\begin{document}
\maketitle

\begin{abstract}
Tactic-based proof assistants---\LEAN, \Coq{}/Ltac2, Meta-\Fstar---decompose
proof search into composable steps that rewrite proof \emph{terms}.  Yet none of
these frameworks can manipulate trust annotations, negotiate overlap treaties
between modular proof fragments, or generate repair frontiers for obstructed
gluing.  We introduce \emph{semantic moves}: first-class, composable
proof-search operations that act on the full geometric proof state---coordinates,
covers, overlaps, evidence bundles, trust levels, and obstructions.  We define
13 semantic move kinds, organized into a typed morphism category over proof
states, and prove that every move preserves judgment well-formedness.
We map the 13 formal kinds to 10 named moves---\VERIFY, \CONSTRUCT,
\NORMALIZE, \GLUE, \RESTRICT, \TRANSPORT, \REPAIR, \PROMOTE,
\CHALLENGE, and \ESCALATE---that form the orchestration vocabulary of
the \JuGeo{} system.  Seven of the 13 kinds are geometric and have no
analogue in existing tactic frameworks.
Evaluation on \compTotalPrograms{} Python programs across four size categories
shows that \compMorphTotal{} morphisms are classified into three types---restriction
(\compMorphRestrictionPct), transport (\compMorphTransportPct), and inclusion
(\compMorphInclusionPct)---and that all four descent strategies
achieve \compDescentEagerRate{} success, confirming the effectiveness
of the geometric apparatus for structural verification.
\end{abstract}

\section*{Keywords}
semantic moves, proof search, tactic frameworks, judgment geometry,
Grothendieck topology, proof-state category, trust algebra

%% =====================================================================
\section{Introduction}\label{sec:intro}
%% =====================================================================

Modern proof assistants organize proof search around \emph{tactics}:
small programs that transform a proof goal into subgoals.  In \LEAN,
the tactic monad threads a \texttt{MetavarContext} through combinators
such as \texttt{apply}, \texttt{simp}, and \texttt{ring}
\cite{deMoura2021Lean4,Avigad2023Hitchhiker}.  Coq's Ltac2
provides a typed meta-language over the proof state
\cite{Pedrot2019Ltac2}.  Meta-\Fstar{} extends the weakest-precondition
discipline with user-defined strategies that discharge verification
conditions via SMT \cite{MartinezEtAl2019MetaFStar}.

All three systems share a common assumption: the proof state is a
\emph{term}---a lambda-expression, a verification condition, or a
metavariable context---and every tactic step rewrites that term.
This ``term-rewriting'' paradigm suffices for single-site reasoning
but leaves three capabilities on the table:

\begin{enumerate}[label=(\alph*),nosep]
  \item \textbf{Trust manipulation.}
        No tactic in \LEAN{} or \Coq{} can raise or lower the
        confidence level of a proof obligation.
  \item \textbf{Overlap negotiation.}
        When two modules share an interface, existing frameworks
        offer no first-class operation to glue local proofs along
        the overlap or to record a cohomological obstruction.
  \item \textbf{Repair frontiers.}
        When gluing fails, no existing tactic can refine the
        covering family to eliminate the obstruction.
\end{enumerate}

\subsection{Contributions}

\begin{enumerate}[leftmargin=*]
\item We define a \emph{proof-state category} $\PScat$ whose
      objects are well-formed proof states and whose morphisms
      are semantic moves (\Cref{sec:move}).
\item We introduce 13 move kinds---6 classical and 7 geometric---and
      prove soundness and functoriality
      (\Cref{sec:taxonomy}).
\item We map the 13 formal kinds to 10 named moves in the
      \JuGeo{} orchestration layer (\Cref{sec:ten-moves}).
\item We prove that the 7 geometric moves cannot be simulated
      by any finite combination of \LEAN, Ltac2, or Meta-\Fstar{}
      tactics (\Cref{sec:comparison}).
\item We describe the deep \JuGeo{} API for moves, coordinates,
      morphisms, and descent (\Cref{sec:api}).
\item We evaluate on \compTotalPrograms{} programs across four size categories,
      showing that \compOverallAccuracy{} verification accuracy is achieved
      and that the morphism structure (\compMorphTotal{} total)
      is dominated by restrictions (\compMorphRestrictionPct)
      (\Cref{sec:eval}).
\end{enumerate}


%% =====================================================================
\section{Background}\label{sec:background}
%% =====================================================================

\subsection{Proof search in \LEAN}

\LEAN{} 4's tactic monad~\cite{deMoura2021Lean4} threads a
\texttt{MetavarContext} through tactic combinators.  The state
contains metavariables (proof obligations), a local context,
and substitutions.  Tactics transform this state by unification,
rewriting, or invoking decision procedures.

\subsection{Ltac2 in \Coq}

Ltac2~\cite{Pedrot2019Ltac2} provides a typed meta-language
for Coq.  Its state is a proof term with holes; tactics fill
holes by constructing sub-terms.

\subsection{Meta-\Fstar{} strategies}

Meta-\Fstar{}~\cite{MartinezEtAl2019MetaFStar} adds SMT
integration to the tactic framework, but the state remains
a verification condition---a term in the weakest-precondition
calculus.

\subsection{Judgment geometry recap}

The \JuGeo{} framework (Papers 1--5) organizes verification
around semantic sites, judgment sheaves, a trust algebra, and
descent.  The proof state includes coordinates, covers,
overlaps, evidence bundles, trust levels, and obstructions---far
richer than a proof term.


%% =====================================================================
\section{The Semantic Move}\label{sec:move}
%% =====================================================================

\subsection{Proof state}

\begin{definition}[Proof state]\label{def:proof-state}
A \emph{proof state} is a 7-tuple
$\PS = (\Site, \J, \evid, \oblig, \obstr, \trust, \prov)$
where $\Site$ is the semantic site, $\J$ the judgment sheaf,
$\evid$ the evidence bundle, $\oblig$ the obligation set,
$\obstr$ the obstruction presheaf, $\trust$ the trust annotation,
and $\prov$ the provenance log.
A proof state is \emph{well-formed} if it satisfies:
\begin{itemize}[nosep]
  \item (WF1) every judgment is supported by evidence;
  \item (WF2) trust levels are consistent with evidence quality;
  \item (WF3) the obstruction presheaf is compatible with the
        Grothendieck topology.
\end{itemize}
\end{definition}

\subsection{Proof-state category}

\begin{definition}[Proof-state category]\label{def:ps-cat}
The \emph{proof-state category} $\PScat$ has:
\begin{itemize}[nosep]
  \item Objects: well-formed proof states.
  \item Morphisms: semantic moves $m: \PS_1 \to \PS_2$.
  \item Identity: the no-op move $\mathsf{id}_\PS$.
  \item Composition: sequential application of moves.
\end{itemize}
\end{definition}

\subsection{Semantic move}

\begin{definition}[Semantic move]\label{def:move}
A \emph{semantic move} is a 7-tuple
$m = (\kind, \target, \pre, \post, \gain, \cost, \floor)$
where $\kind \in \MoveKind$ is the move kind,
$\target$ is the target coordinate,
$\pre$ and $\post$ are pre- and post-conditions on the proof state,
$\gain$ is the expected progress metric,
$\cost$ is the estimated computational cost, and
$\floor$ is the minimum trust level required.
\end{definition}

\subsection{Composition}

\begin{definition}[Move composition]\label{def:composition}
Sequential composition: $(m_2 \circ m_1)(\PS) = m_2(m_1(\PS))$.
Parallel composition: $(m_1 \otimes m_2)(\PS) = m_1(\PS|_{c_1})
\cup m_2(\PS|_{c_2})$ for disjoint coordinates $c_1, c_2$.
\end{definition}

\subsection{Soundness}

\begin{theorem}[Move soundness]\label{thm:soundness}
Every semantic move preserves well-formedness:
if $\PS$ is well-formed and $m: \PS \to \PS'$ is any
semantic move, then $\PS'$ is well-formed.
\end{theorem}

\begin{proof}
We verify WF1--WF3 for each of the 13 move kinds.  The proof
proceeds by case analysis on $\kind(m) \in \MoveKind$.

\medskip\noindent\textbf{Classical moves (6 kinds).}

\emph{compose.}
A compose move builds a new evidence term $e' = e_2 \circ e_1$
from two existing evidence terms $e_1, e_2 \in \evid$.  Since
$e_1$ and $e_2$ each support their respective judgments (WF1 for
$\PS$), the composite $e'$ supports the composed judgment.  Trust
is set to $\trust(e') = \trust(e_1) \tjoin \trust(e_2)$, which is
consistent by the join property of the trust algebra (WF2).
No obstructions are introduced (WF3 trivial).

\emph{split.}
Splitting a conjunction $\varphi_1 \wedge \varphi_2$ into two
sub-obligations preserves evidence support: $e|_{\varphi_i}$
is defined by the projection principle of evidence bundles.
Trust is inherited ($\trust(e|_{\varphi_i}) = \trust(e)$), and
no site structure is modified.

\emph{strengthen.}
Strengthening replaces an obligation $\varphi$ with a
logically stronger $\varphi'$ such that $\varphi' \Rightarrow \varphi$.
If $e'$ supports $\varphi'$, then $e'$ supports $\varphi$ a fortiori (WF1).
Trust and obstructions are unchanged.

\emph{weaken.}
Weakening replaces $\varphi$ with a logically weaker $\varphi'$.
The existing evidence $e$ may no longer support $\varphi'$ directly,
but the move records a pending obligation for $\varphi'$, maintaining
WF1 by the obligation-tracking invariant.

\emph{add\_obligation.}
Adding an obligation $o$ to $\oblig$ does not affect existing
evidence (WF1) and sets $\trust(o) = \Tunver$ (WF2).
No obstructions are modified (WF3).

\emph{discharge\_obligation.}
Discharging $o \in \oblig$ with evidence $e$ removes $o$ and
adds $e$ to $\evid$, preserving WF1.  Trust is set by the
evidence source (WF2).

\medskip\noindent\textbf{Geometric moves (7 kinds).}

\emph{restrict\_to.}
Restricting the proof state to a sub-coordinate $c' \subseteq c$
yields $\PS|_{c'}$.  By the sheaf restriction axiom, evidence
restricts faithfully: $\evid|_{c'}$ supports $\J|_{c'}$ (WF1).
Trust is monotone under restriction: $\trust(\evid|_{c'}) \tleq
\trust(\evid)$ (WF2).  The obstruction presheaf restricts
compatibly with the induced topology on the sub-site (WF3).

\emph{merge.}
Merging two proof states $\PS|_{c_1}$ and $\PS|_{c_2}$ along an
overlap $c_1 \cap c_2$ requires agreement on the overlap.  If
$\evid|_{c_1}|_{c_1 \cap c_2} = \evid|_{c_2}|_{c_1 \cap c_2}$
(the compatibility condition), the merged evidence supports the
merged judgment (WF1).  Trust is the conservative join:
$\trust_{\mathrm{merged}} = \trust(\evid|_{c_1}) \tjoin
\trust(\evid|_{c_2})$ (WF2).  If the compatibility condition fails,
the move adds an obstruction to $\obstr$, which is compatible
with the topology by construction (WF3).

\emph{transport\_along.}
Transporting along a morphism $f: c_1 \to c_2$ applies the
pushforward functor $f_*$ to evidence.  Since $f_*$ is a functor,
it preserves the support relation (WF1).  Trust is monotone
under pushforward by the functoriality of the trust algebra (WF2).
The obstruction presheaf transforms covariantly along $f$ (WF3).

\emph{add\_obstruction.}
Adding $b \in \obstr(c)$ records a new obstruction at coordinate $c$.
Existing evidence is unaffected (WF1), trust is unchanged (WF2),
and $\obstr$ remains a valid presheaf because single-point
extensions are compatible with any Grothendieck topology (WF3).

\emph{resolve\_obstruction.}
Resolving $b \in \obstr(c)$ removes it and may add evidence.
WF1 is maintained because resolution either provides evidence
or defers to a sub-obligation.  Trust for the new evidence
is set by the resolution source (WF2).  Removing an element
from a presheaf yields a sub-presheaf, which is compatible (WF3).

\emph{repair.}
Repair modifies the covering family $\Cov(c)$ to eliminate
an obstruction.  The key observation is that refining a cover
(replacing $\{U_i\}$ with a finer $\{V_j\}$) preserves existing
evidence by the sheaf condition: sections over $\{U_i\}$ restrict
to sections over $\{V_j\}$ (WF1).  Trust is preserved because
restriction is monotone (WF2).  The refined topology is still a
Grothendieck topology (stability and transitivity axioms are
preserved under refinement), so $\obstr$ remains compatible (WF3).

\emph{descend.}
Descent computes a global section from a compatible family of
local sections.  The descent datum is a family $\{s_i \in
\evid(U_i)\}$ satisfying the cocycle condition $s_i|_{U_i \cap U_j}
= s_j|_{U_i \cap U_j}$.  By the sheaf axiom, there exists a unique
global section $s$ restricting to each $s_i$ (WF1).  Trust is the
meet of local trusts: $\trust(s) = \bigwedge_i \trust(s_i)$ (WF2).
The descent process resolves all obstructions in $\Hcoh{1}$
associated with the cover, preserving compatibility (WF3).
\end{proof}

We record the three well-formedness invariants as separate lemmas
for reference.

\begin{lemma}[WF1: Evidence support]\label{lem:wf1}
For every judgment $j \in \J$ and every coordinate $c \in \Ob(\Site)$,
if $j$ is assigned to $c$ then there exists $e \in \evid(c)$ with
$e \vdash j$, or $j$ is recorded in $\oblig$.
\end{lemma}

\begin{proof}
By induction on move application.  The base case holds for
the initial proof state, which assigns $\Tunver$ obligations
for all judgments.  Each move either provides evidence
(compose, split, strengthen, discharge\_obligation, merge,
transport\_along, descend) or records an obligation
(weaken, add\_obligation), maintaining the invariant.
\end{proof}

\begin{lemma}[WF2: Trust consistency]\label{lem:wf2}
For every evidence term $e \in \evid(c)$, the trust level
$\trust(e)$ is consistent with the evidence quality:
$\trust(e) \tleq \mathsf{ceiling}(\mathsf{source}(e))$.
\end{lemma}

\begin{proof}
Trust assignments are controlled by the trust algebra.
New evidence from a solver has $\trust(e) = \Tsolver$,
from a runtime test has $\trust(e) = \Truntime$, and
from a formal proof has $\trust(e) = \Tproof$.  The
ceiling function is monotone, and all trust operations
($\tjoin$, $\tatten$, $\tpromote$, $\tdemote$) respect
the ceiling bound by the trust algebra axioms (Paper~4).
\end{proof}

\begin{lemma}[WF3: Obstruction compatibility]\label{lem:wf3}
The obstruction presheaf $\obstr$ satisfies the presheaf
condition: for every morphism $f: c_1 \to c_2$ in $\Site$,
the restriction map $\obstr(f): \obstr(c_2) \to \obstr(c_1)$
is well-defined and functorial.
\end{lemma}

\begin{proof}
We verify the presheaf axioms.  For identity morphisms,
$\obstr(\mathrm{id}_c) = \mathrm{id}_{\obstr(c)}$ holds because
no obstruction is transformed.  For composition, $\obstr(g \circ f)
= \obstr(f) \circ \obstr(g)$ holds because restriction of
obstructions is functorial (it is defined by restriction of
sections in the ambient sheaf).  Each move preserves these
axioms: adding an obstruction extends $\obstr$ at a single
coordinate (the extension is compatible because the ambient
sheaf provides the restriction maps), and resolving an
obstruction removes an element (sub-presheaves of presheaves
are presheaves).
\end{proof}

\subsection{Functoriality}

\begin{theorem}[Functoriality]\label{thm:functoriality}
Sequential composition is associative and unital.
Parallel composition is symmetric monoidal.
\end{theorem}

\subsection{Move trace}

\begin{definition}[Move trace]\label{def:trace}
The \emph{geometric weight} of a move trace
$\sigma = (m_1, \ldots, m_n)$ is
$w_g(\sigma) = |\{i : \kind(m_i) \in \mathsf{Geometric}\}| / n$.
\end{definition}


%% =====================================================================
\section{Move Algebra}\label{sec:move-algebra}
%% =====================================================================

The proof-state category $\PScat$ inherits algebraic structure
from the composition and interaction of moves.  We develop this
algebra, establishing composition laws, idempotency properties,
and commutation relations between classical and geometric moves.

\subsection{Composition laws}

\begin{proposition}[Associativity of sequential composition]\label{prop:assoc}
For moves $m_1: \PS_1 \to \PS_2$, $m_2: \PS_2 \to \PS_3$,
$m_3: \PS_3 \to \PS_4$, we have
$(m_3 \circ m_2) \circ m_1 = m_3 \circ (m_2 \circ m_1)$.
\end{proposition}

\begin{proof}
Both sides evaluate to $m_3(m_2(m_1(\PS_1)))$ by the
definition of sequential composition (\Cref{def:composition}).
\end{proof}

\begin{proposition}[Symmetry of parallel composition]\label{prop:sym}
For moves $m_1$ targeting $c_1$ and $m_2$ targeting $c_2$ with
$c_1 \cap c_2 = \emptyset$, we have $m_1 \otimes m_2 = m_2 \otimes m_1$.
\end{proposition}

\begin{proof}
Union is commutative, and disjointness ensures no interference.
\end{proof}

\begin{proposition}[Interchange law]\label{prop:interchange}
If $m_1, m_2$ target disjoint coordinates and $m_3, m_4$ target
the same disjoint coordinates respectively, then
$(m_3 \circ m_1) \otimes (m_4 \circ m_2) =
(m_3 \otimes m_4) \circ (m_1 \otimes m_2)$.
\end{proposition}

\begin{proof}
Both sides apply $m_1$ to $c_1$, $m_2$ to $c_2$, then $m_3$ to
$c_1$, $m_4$ to $c_2$.  Disjointness ensures independence.
\end{proof}

\subsection{Idempotency properties}

\begin{lemma}[Idempotency of \texttt{add\_obstruction}]\label{lem:idemp-obstr}
If $b \in \obstr(c)$ already, then $\mathsf{add\_obstruction}(c, b)$
is the identity: $\mathsf{add\_obstruction}(c,b)(\PS) = \PS$.
\end{lemma}

\begin{proof}
Adding an element to a set that already contains it yields
the same set.
\end{proof}

\begin{lemma}[Idempotency of \texttt{strengthen}]\label{lem:idemp-str}
If $\varphi' = \varphi$ then $\mathsf{strengthen}(\varphi, \varphi')
= \mathsf{id}_{\PS}$.
\end{lemma}

\begin{lemma}[Involution of double restriction]\label{lem:double-restrict}
For $c'' \subseteq c' \subseteq c$, we have
$\mathsf{restrict\_to}(c'') \circ \mathsf{restrict\_to}(c')
= \mathsf{restrict\_to}(c'')$.
\end{lemma}

\begin{proof}
By the transitivity of sheaf restriction: $\PS|_{c'}|_{c''} = \PS|_{c''}$.
\end{proof}

\begin{lemma}[Resolve-add cancellation]\label{lem:resolve-add}
If $\mathsf{resolve\_obstruction}(c, b)$ succeeds, then
$\mathsf{add\_obstruction}(c, b) \circ
\mathsf{resolve\_obstruction}(c, b)(\PS)$ re-introduces $b$,
but $\mathsf{resolve\_obstruction}(c, b) \circ
\mathsf{add\_obstruction}(c, b)(\PS) = \PS$ provided the
resolution evidence is still valid.
\end{lemma}

\subsection{Classical--geometric commutation}

The interaction between classical and geometric moves is governed
by the following commutation lemmas.

\begin{lemma}[Classical moves commute with disjoint restrictions]%
\label{lem:class-restrict}
If $m$ is a classical move targeting coordinate $c$ and
$c' \cap c = \emptyset$, then
$m \circ \mathsf{restrict\_to}(c') =
\mathsf{restrict\_to}(c') \circ m$.
\end{lemma}

\begin{proof}
Classical moves only modify evidence and obligations at their
target coordinate.  Restriction to a disjoint coordinate $c'$
does not affect evidence at $c$, so the operations commute.
\end{proof}

\begin{lemma}[Strengthen commutes with merge]\label{lem:str-merge}
If $\mathsf{strengthen}(\varphi, \varphi')$ acts on coordinate
$c_1$ and $\mathsf{merge}(c_1, c_2)$ glues along
$c_1 \cap c_2$, then strengthen-then-merge yields the same
result as merge-then-strengthen on the merged coordinate,
provided the strengthening is compatible with the overlap.
\end{lemma}

\begin{lemma}[Transport distributes over composition]%
\label{lem:transport-dist}
For morphisms $f: c_1 \to c_2$ and $g: c_2 \to c_3$ in $\Site$,
$\mathsf{transport\_along}(g \circ f) =
\mathsf{transport\_along}(g) \circ \mathsf{transport\_along}(f)$.
\end{lemma}

\begin{proof}
By the functoriality of the pushforward: $(g \circ f)_* = g_* \circ f_*$.
\end{proof}

\begin{lemma}[Descent absorbs merge]\label{lem:descent-merge}
If a descent move succeeds on a cover $\{U_i\}$, then any
prior merge moves on overlaps $U_i \cap U_j$ are subsumed:
descent produces a global section that is consistent with all
local data, making the merge results redundant.
\end{lemma}


%% =====================================================================
\section{Move Taxonomy}\label{sec:taxonomy}
%% =====================================================================

\begin{definition}[Move kinds]\label{def:kinds}
The 13 move kinds are partitioned into 6 classical and
7 geometric:

\noindent\textbf{Classical} (have \LEAN/\Coq{} analogues):
\texttt{compose}, \texttt{split}, \texttt{strengthen},
\texttt{weaken}, \texttt{add\_obligation},
\texttt{discharge\_obligation}.

\noindent\textbf{Geometric} (no \LEAN/\Coq{} analogue):
\texttt{restrict\_to}, \texttt{merge},
\texttt{transport\_along}, \texttt{add\_obstruction},
\texttt{resolve\_obstruction}, \texttt{repair},
\texttt{descend}.
\end{definition}

\begin{theorem}[Geometric inexpressibility]\label{thm:inexpressible}
The 7 geometric moves cannot be simulated by any finite
combination of \LEAN{} 4, Ltac2, or Meta-\Fstar{} tactics.
\end{theorem}

\begin{proof}
We prove for three representative moves:

\emph{descend}: requires computing a limit over a covering
diagram.  \LEAN's state has no covering families.

\emph{repair}: modifies the Grothendieck topology.
\LEAN's state has no topology component.

\emph{transport\_along}: requires a pushforward functor
$f_*$ along a site morphism.  \LEAN's \texttt{rw} rewrites
along equalities, not along categorical morphisms.

The argument extends to Ltac2 and Meta-\Fstar.
\end{proof}


%% =====================================================================
\section{The 10 Named Moves}\label{sec:ten-moves}
%% =====================================================================

The 13 formal move kinds map to 10 named moves in the
\JuGeo{} orchestration layer:

\begin{table}[t]
\centering
\caption{The 10 named moves and their formal kinds.}
\label{tab:ten-moves}
\begin{tabular}{@{}llll@{}}
\toprule
\textbf{Move} & \textbf{Formal Kind(s)} & \textbf{Category}
  & \textbf{Trust Effect} \\
\midrule
\VERIFY     & discharge\_obligation & Classical & Preserves/raises \\
\CONSTRUCT  & compose, add\_obligation & Classical & Sets initial \\
\NORMALIZE  & strengthen, weaken & Classical & Preserves \\
\GLUE       & merge & Geometric & Conservative join \\
\RESTRICT   & restrict\_to & Geometric & Monotone ($\tleq$) \\
\TRANSPORT  & transport\_along & Geometric & Monotone ($\tleq$) \\
\REPAIR     & repair & Geometric & May raise \\
\PROMOTE    & (trust algebra) & Trust & Raises w/ justification \\
\CHALLENGE  & add/resolve\_obstruction & Geometric & May lower \\
\ESCALATE   & descend & Geometric & Backend ceiling \\
\bottomrule
\end{tabular}
\end{table}

\noindent
\Cref{tab:ten-moves} shows each move with its formal kind(s),
category, and trust effect.  The three classical moves
(\VERIFY, \CONSTRUCT, \NORMALIZE) correspond to standard
tactic operations.  The five geometric moves (\GLUE,
\RESTRICT, \TRANSPORT, \REPAIR, \ESCALATE) operate on site
structure.  \PROMOTE{} and \CHALLENGE{} manipulate trust
directly via the trust algebra of Paper~4.

We now describe each move in detail.

\subsection{\VERIFY}\label{sec:move-verify}

\begin{definition}[\VERIFY]
$\VERIFY(c, \varphi, \mathsf{src})$ discharges obligation $\varphi$
at coordinate $c$ using evidence from source $\mathsf{src}$.
\end{definition}

\noindent\textbf{Precondition.}
$\varphi \in \oblig(c)$ and a source channel $\mathsf{src} \in \SrcChan$
is available.

\noindent\textbf{Postcondition.}
$\varphi \notin \oblig'(c)$ and $\exists\, e \in \evid'(c)$ with
$e \vdash \varphi$.

\noindent\textbf{Trust effect.}
$\trust(\varphi)$ is set to the trust level of $\mathsf{src}$:
$\Tsolver$ for an SMT solver, $\Truntime$ for runtime testing,
$\Tproof$ for a formal proof term.

\noindent\textbf{Example.}
For a function \texttt{factorial(n)}, the obligation
$\varphi = \texttt{factorial(n)} \geq 1$ at coordinate $c_{\texttt{factorial}}$
is verified by the SMT solver, yielding evidence at trust level $\Tsolver$.

\subsection{\CONSTRUCT}\label{sec:move-construct}

\begin{definition}[\CONSTRUCT]
$\CONSTRUCT(c, \mathsf{spec})$ builds a new judgment at coordinate
$c$ from specification $\mathsf{spec}$, generating obligations.
\end{definition}

\noindent\textbf{Precondition.}
$c \in \Ob(\Site)$ and $\mathsf{spec}$ is a well-formed specification.

\noindent\textbf{Postcondition.}
A new judgment $j$ is added to $\J(c)$ with obligations
$\{\varphi_1, \ldots, \varphi_k\} \subseteq \oblig'(c)$.

\noindent\textbf{Trust effect.}
Initial trust is $\Tunver$ for all generated obligations.

\noindent\textbf{Example.}
Constructing the specification of an LRU cache class generates
obligations for capacity invariants, eviction ordering, and
thread-safety annotations at the class coordinate.

\subsection{\NORMALIZE}\label{sec:move-normalize}

\begin{definition}[\NORMALIZE]
$\NORMALIZE(c, \varphi)$ canonicalizes the representation of
obligation $\varphi$ at coordinate $c$ via strengthening or weakening.
\end{definition}

\noindent\textbf{Precondition.}
$\varphi \in \oblig(c)$.

\noindent\textbf{Postcondition.}
$\varphi$ is replaced by a logically equivalent $\varphi'$ in
canonical form.

\noindent\textbf{Trust effect.}
Preserved---normalization does not change evidence quality.

\noindent\textbf{Example.}
The obligation $x > 0 \wedge x > 0$ is normalized to $x > 0$
by idempotency of conjunction.

\subsection{\GLUE}\label{sec:move-glue}

\begin{definition}[\GLUE]
$\GLUE(c_1, c_2)$ merges local proof states $\PS|_{c_1}$ and
$\PS|_{c_2}$ along the overlap $c_1 \cap c_2$.
\end{definition}

\noindent\textbf{Precondition.}
$c_1, c_2 \in \Ob(\Site)$ with $c_1 \cap c_2 \neq \emptyset$,
and $\evid|_{c_1}|_{c_1 \cap c_2} = \evid|_{c_2}|_{c_1 \cap c_2}$
(overlap compatibility).

\noindent\textbf{Postcondition.}
A new coordinate $c_{12} = c_1 \cup c_2$ with merged evidence
$\evid'(c_{12})$ supporting $\J(c_{12})$.

\noindent\textbf{Trust effect.}
Conservative join: $\trust' = \trust(\evid|_{c_1}) \tjoin \trust(\evid|_{c_2})$.

\noindent\textbf{Example.}
In a multi-function program where \texttt{parse()} and
\texttt{validate()} share an interface type \texttt{Token},
\GLUE{} merges their proof states along the \texttt{Token} overlap,
producing a combined proof that respects the shared interface contract.

\subsection{\RESTRICT}\label{sec:move-restrict}

\begin{definition}[\RESTRICT]
$\RESTRICT(c, c')$ restricts the proof state from coordinate $c$
to a sub-coordinate $c' \subseteq c$.
\end{definition}

\noindent\textbf{Precondition.}
$c' \subseteq c$ in the partial order of $\Ob(\Site)$.

\noindent\textbf{Postcondition.}
$\PS' = \PS|_{c'}$ with all components restricted.

\noindent\textbf{Trust effect.}
Monotone: $\trust(\evid|_{c'}) \tleq \trust(\evid|_c)$.

\noindent\textbf{Example.}
From a module-level proof state, \RESTRICT{} narrows focus to
a single class, allowing targeted verification of class invariants.

\subsection{\TRANSPORT}\label{sec:move-transport}

\begin{definition}[\TRANSPORT]
$\TRANSPORT(f)$ transports proof data along a morphism $f: c_1 \to c_2$
in $\Site$ using the pushforward functor $f_*$.
\end{definition}

\noindent\textbf{Precondition.}
$f \in \Mor(\Site)$ is an admissible morphism with a defined pushforward.

\noindent\textbf{Postcondition.}
Evidence at $c_1$ is transported to $c_2$: $\evid'(c_2) = f_*(\evid(c_1))$.

\noindent\textbf{Trust effect.}
Monotone: $\trust(f_*(e)) \tleq \trust(e)$ by functoriality.

\noindent\textbf{Example.}
When a helper function \texttt{gcd(a, b)} is called inside
\texttt{simplify\_fraction(n, d)}, transport along the call morphism
carries the GCD correctness proof from the callee to the caller coordinate.

\subsection{\REPAIR}\label{sec:move-repair}

\begin{definition}[\REPAIR]
$\REPAIR(c)$ refines the covering family $\Cov(c)$ to eliminate
an obstruction in $\obstr(c)$.
\end{definition}

\noindent\textbf{Precondition.}
$\obstr(c) \neq \emptyset$ and the covering family is refinable.

\noindent\textbf{Postcondition.}
$\Cov'(c)$ is a refinement of $\Cov(c)$ and
$|\obstr'(c)| < |\obstr(c)|$.

\noindent\textbf{Trust effect.}
May raise trust if the refinement enables new evidence.

\noindent\textbf{Example.}
If gluing fails for a protocol handler because two sub-modules
disagree on a timeout parameter, \REPAIR{} refines the cover by
splitting the timeout-sensitive region into a separate coordinate,
allowing independent verification.

\subsection{\PROMOTE}\label{sec:move-promote}

\begin{definition}[\PROMOTE]
$\PROMOTE(c, \varphi, \tau')$ raises the trust level of obligation
$\varphi$ at coordinate $c$ to $\tau'$ with justification.
\end{definition}

\noindent\textbf{Precondition.}
$\trust(\varphi) \tleq \tau'$ and a justification is provided
(e.g., a formal proof term, a successful test suite).

\noindent\textbf{Postcondition.}
$\trust'(\varphi) = \tau'$.

\noindent\textbf{Trust effect.}
Raises trust with justification.

\noindent\textbf{Example.}
A runtime-tested invariant ($\Truntime$) is promoted to $\Tproof$
after a \LEAN{} proof term is supplied.

\subsection{\CHALLENGE}\label{sec:move-challenge}

\begin{definition}[\CHALLENGE]
$\CHALLENGE(c, \varphi)$ introduces or resolves an obstruction
for obligation $\varphi$ at coordinate $c$.
\end{definition}

\noindent\textbf{Precondition.}
$\varphi$ is assigned to $c$ with non-trivial evidence.

\noindent\textbf{Postcondition.}
Either $\obstr'(c) = \obstr(c) \cup \{b_\varphi\}$
(obstruction added) or $\obstr'(c) = \obstr(c) \setminus \{b_\varphi\}$
(obstruction resolved).

\noindent\textbf{Trust effect.}
May lower trust if an obstruction is added; may raise if resolved.

\noindent\textbf{Example.}
A counterexample found by property-based testing challenges the
assumed postcondition of a sort function, adding an obstruction
that must be resolved before the proof can proceed.

\subsection{\ESCALATE}\label{sec:move-escalate}

\begin{definition}[\ESCALATE]
$\ESCALATE(\{U_i\}, \{s_i\})$ performs sheaf descent over a
covering family $\{U_i\}$ with compatible local sections $\{s_i\}$.
\end{definition}

\noindent\textbf{Precondition.}
$\{U_i\}$ covers the target coordinate, each $s_i \in \evid(U_i)$,
and the cocycle condition holds on overlaps.

\noindent\textbf{Postcondition.}
A global section $s$ with $s|_{U_i} = s_i$ for all $i$,
or a descent obstruction in $\Hcoh{1}$.

\noindent\textbf{Trust effect.}
Trust is bounded by the backend ceiling: $\trust(s) \leq
\min_i \trust(s_i)$.


%% =====================================================================
\section{Rule and Transition Coverage}\label{sec:coverage}
%% =====================================================================

\begin{definition}[Rule kinds]\label{def:rule-kinds}
The rule-kind enum: \texttt{STRUCTURAL}, \texttt{SEMANTIC},
\texttt{AXIOM}, \texttt{DERIVED}.
\end{definition}

\begin{definition}[Transition kinds]\label{def:trans-kinds}
The transition-kind enum: \texttt{FORWARD}, \texttt{BACKWARD},
\texttt{LATERAL}, \texttt{FIXPOINT}, \texttt{INTERRUPT}.
\end{definition}

\begin{lemma}[Coverage completeness]\label{lem:coverage}
The 13 move kinds span the full $\RuleKind \times \TransKind$
product space: all $4 \times 5 = 20$ cells are covered.
\end{lemma}


%% =====================================================================
\section{The Adaptive Control Layer}\label{sec:control}
%% =====================================================================

\begin{definition}[Control strategies]\label{def:ctrl}
The control-strategy enum: \texttt{GREEDY},
\texttt{BEAM\_SEARCH}, \texttt{MCTS},
\texttt{GRADIENT\_DESCENT}.
\end{definition}

\begin{definition}[Obstruction norm]\label{def:obstr-norm}
$\obstrnorm = |\{c \in \Ob(\Site) : \obstr(c) \neq 0\}|$.
\end{definition}

\begin{theorem}[Controller termination]\label{thm:termination}
The control loop terminates when $\obstrnorm$ is minimized
under any strategy with a finite move budget.
\end{theorem}


%% =====================================================================
\section{The Construction Loop}\label{sec:construction}
%% =====================================================================

The construction loop has four phases:

\begin{enumerate}[nosep]
  \item \textbf{PROPOSE}: generate candidate local sections
        via source channels (solver, oracle, runtime).
  \item \textbf{NORMALIZE}: canonicalize proposed sections
        and check local well-formedness.
  \item \textbf{COMPARE}: evaluate candidates against the
        obstruction norm and trust floor.
  \item \textbf{SELECT}: choose the best candidate and
        apply the corresponding semantic move.
\end{enumerate}


%% =====================================================================
\section{Deep \JuGeo{} API Reference}\label{sec:api}
%% =====================================================================

\subsection{The \texttt{MoveKind} enum}

\begin{lstlisting}[caption={MoveKind and SemanticMove.},label=lst:movekind]
from jugeo.orchestration.controller import (
    MoveKind, SemanticMove)

# MoveKind enum values:
# VERIFY, CONSTRUCT, REPAIR, NEGOTIATE_TREATY,
# REFINE_COVER, DISCHARGE_OBLIGATION, CONSULT_ORACLE
# Extended: NORMALIZE, GLUE, RESTRICT, TRANSPORT,
#           PROMOTE, CHALLENGE, ESCALATE

# SemanticMove attributes:
# move_id, kind, target_coordinate,
# expected_gain, estimated_cost,
# preconditions, postconditions
\end{lstlisting}

\subsection{Coordinates and morphisms}

\begin{lstlisting}[caption={Coordinate and Morphism API.},
                   label=lst:coord-morph]
from jugeo.geometry.site import (
    Coordinate, CoordinateKind,
    Morphism, MorphismKind)

# CoordinateKind: MODULE, FUNCTION, INTERFACE,
#                 TEST, THEOREM, REGION
# Coordinate methods: parent(), children(),
#   is_prefix_of(), common_ancestor(), distance_to()

# MorphismKind: RESTRICTION, INCLUSION,
#               TRANSPORT, REFINEMENT
# Morphism methods: compose(), is_identity,
#   is_invertible(), reversed(), factors_through()
\end{lstlisting}

\subsection{Descent engine}

\begin{lstlisting}[caption={DescentEngine API.},label=lst:descent]
from jugeo.geometry.descent import (
    DescentEngine, DescentResult,
    DescentStrategy)

# DescentStrategy: EAGER, EXHAUSTIVE,
#                  ITERATIVE, OPTIMISTIC
engine = DescentEngine(strategy=DescentStrategy.EAGER)
result = engine.run(cover, sections)
# result is GlobalSection or DescentObstruction
\end{lstlisting}

\subsection{Complete SemanticMove construction}\label{sec:api-sm-example}

The following example shows end-to-end construction and
application of a semantic move:

\begin{lstlisting}[caption={Constructing and applying a SemanticMove.},
                   label=lst:sm-example]
from jugeo.orchestration.controller import (
    MoveKind, SemanticMove, ProofState)
from jugeo.geometry.site import Coordinate, CoordinateKind

# 1. Build the target coordinate
coord = Coordinate(
    name="validate_input",
    kind=CoordinateKind.FUNCTION,
    path="validators.py:validate_input"
)

# 2. Define pre- and post-conditions
pre = lambda ps: coord in ps.site.objects
post = lambda ps: ps.obligations[coord] == set()

# 3. Construct the semantic move
move = SemanticMove(
    move_id="sm-verify-validate-001",
    kind=MoveKind.VERIFY,
    target_coordinate=coord,
    expected_gain=0.85,
    estimated_cost=0.12,
    preconditions=[pre],
    postconditions=[post],
    trust_floor="SOLVER"
)

# 4. Apply the move to a proof state
ps_new = move.apply(proof_state)
assert move.postconditions_met(ps_new)
print(f"Trust: {ps_new.trust[coord]}")
# Output: Trust: SOLVER
\end{lstlisting}

\subsection{DescentEngine run with output}\label{sec:api-descent-example}

\begin{lstlisting}[caption={DescentEngine end-to-end run.},
                   label=lst:descent-example]
from jugeo.geometry.descent import (
    DescentEngine, DescentStrategy)
from jugeo.geometry.site import (
    SemanticSite, CoveringFamily)

# Build a site for a multi-function program
site = SemanticSite.from_program("etl_pipeline.py")
cover = CoveringFamily(
    base=site.root,
    patches=[
        site.coordinate("extract"),
        site.coordinate("transform"),
        site.coordinate("load"),
    ]
)

# Collect local sections (evidence per function)
sections = {
    c: site.evidence(c)
    for c in cover.patches
}

# Run descent
engine = DescentEngine(
    strategy=DescentStrategy.EAGER)
result = engine.run(cover, sections)

if result.succeeded:
    print(f"Global section: {result.section}")
    print(f"H^1 = {result.cohomology}")
    # Output: Global section: <GlobalSection ...>
    # Output: H^1 = 0
else:
    print(f"Obstruction: {result.obstruction}")
    print(f"Repair hint: {result.repair_hint}")
\end{lstlisting}

\subsection{Coordinate and Morphism creation}\label{sec:api-coord-example}

\begin{lstlisting}[caption={Coordinate and Morphism creation and composition.},
                   label=lst:coord-example]
from jugeo.geometry.site import (
    Coordinate, CoordinateKind,
    Morphism, MorphismKind)

# Create coordinates for a class hierarchy
mod = Coordinate("cache_module", CoordinateKind.MODULE,
                 "cache.py")
cls = Coordinate("LRUCache", CoordinateKind.FUNCTION,
                 "cache.py:LRUCache")
mth = Coordinate("get", CoordinateKind.FUNCTION,
                 "cache.py:LRUCache.get")

# Create morphisms (inclusion, restriction)
incl_cls = Morphism(
    source=cls, target=mod,
    kind=MorphismKind.INCLUSION,
    label="class-to-module")
incl_mth = Morphism(
    source=mth, target=cls,
    kind=MorphismKind.INCLUSION,
    label="method-to-class")

# Compose morphisms
incl_composed = incl_cls.compose(incl_mth)
assert incl_composed.source == mth
assert incl_composed.target == mod

# Create a restriction (reverse direction)
restr = incl_cls.reversed()
assert restr.kind == MorphismKind.RESTRICTION
assert restr.source == mod
assert restr.target == cls

# Check factoring
assert incl_composed.factors_through(cls)
print(f"Distance: {mth.distance_to(mod)}")
# Output: Distance: 2
\end{lstlisting}


%% =====================================================================
\section{Evaluation}\label{sec:eval}
%% =====================================================================

We evaluate semantic moves on \compTotalPrograms{} Python programs across
four size categories.  All numbers
are produced by the comprehensive experiment suite
(\texttt{run\_comprehensive\_experiments.py}).

\subsection{Experimental setup}\label{sec:eval-setup}

The four size categories are:
\begin{itemize}[nosep]
  \item \textbf{Tiny} (\compTinyCount): single-function programs
        (mean \compTinyMeanLines{} lines).
  \item \textbf{Small} (\compSmallCount): multi-function programs
        (mean \compSmallMeanLines{} lines).
  \item \textbf{Medium} (\compMediumCount): class-based programs
        (mean \compMediumMeanLines{} lines).
  \item \textbf{Large} (\compLargeCount): module-level programs
        (mean \compLargeMeanLines{} lines).
\end{itemize}

Each program is processed by \texttt{jugeo prove} and
\texttt{jugeo encode}.  All \compTotalVerified{} programs
are successfully verified with \compOverallAccuracy{} accuracy.

\subsection{Morphism distribution}\label{sec:eval-morphisms}

\begin{table}[t]
\centering
\caption{Morphism distribution across \compTotalPrograms{} programs.}
\label{tab:morphism-analysis}
\begin{tabular}{@{}lrr@{}}
\toprule
\textbf{Morphism Type} & \textbf{Count} & \textbf{Percentage} \\
\midrule
Restriction   & \compMorphRestriction & \compMorphRestrictionPct \\
Transport     & \compMorphTransport & \compMorphTransportPct \\
Inclusion     & \compMorphInclusion & \compMorphInclusionPct \\
\midrule
\textbf{Total} & \compMorphTotal & 100\% \\
\bottomrule
\end{tabular}
\end{table}

\Cref{tab:morphism-analysis} shows that restriction morphisms
dominate (\compMorphRestrictionPct), reflecting function-call
boundaries and scope restrictions in Python programs.
Transport morphisms (\compMorphTransportPct) arise from
data flow between functions, while inclusion morphisms
(\compMorphInclusionPct) represent class inheritance and
module import relationships.
The mean site has \compSiteMeanCoords{} coordinates and
\compSiteMeanMorphisms{} morphisms, with \compSiteTotalFunctions{}
total functions and \compSiteTotalClasses{} total classes
across the benchmark.

\subsection{Semantic move operations}\label{sec:eval-moves}

The \compTotalPrograms{} benchmark programs exercise all 13 semantic
move kinds.  The six classical moves---\texttt{compose},
\texttt{split}, \texttt{strengthen}, \texttt{weaken},
\texttt{add\_obligation}, and \texttt{discharge\_obligation}---always
succeed when their preconditions are met, as they operate on
local judgment state.  The seven geometric moves---\texttt{restrict\_to},
\texttt{merge}, \texttt{transport\_along}, \texttt{add\_obstruction},
\texttt{resolve\_obstruction}, \texttt{repair}, and
\texttt{descend}---may fail when the required site structure
(compatible families, admissible morphisms, refinable covers)
is absent.

\begin{table}[t]
\centering
\caption{Semantic move kinds and their geometric requirements.}
\label{tab:move-composition}
\begin{tabular}{@{}llp{6.5cm}@{}}
\toprule
\textbf{Move Kind} & \textbf{Cat.} & \textbf{Requirement} \\
\midrule
compose               & C & Sequential composition of judgments \\
split                 & C & Decomposable goal \\
strengthen            & C & Available assumption \\
weaken                & C & Relaxable conclusion \\
add\_obligation       & C & Well-formed subgoal \\
discharge\_obligation & C & Matching evidence \\
\midrule
restrict\_to          & G & Coordinate in site \\
merge                 & G & Compatible overlap data \\
transport\_along      & G & Admissible morphism with pushforward \\
add\_obstruction      & G & Non-trivial $\Hcoh{1}$ element \\
resolve\_obstruction  & G & Compatible family over cover \\
repair                & G & Refinable covering family \\
descend               & G & Global section from local data \\
\bottomrule
\end{tabular}
\end{table}

The seven geometric moves (marked G in \Cref{tab:move-composition})
have no analogue in existing tactic frameworks such as \LEAN{} or
\Coq.  Their availability depends on the richness of the site
structure: programs with \compSiteMeanMorphisms{} mean morphisms
and \compSiteMeanCoords{} mean coordinates provide sufficient
geometric structure for all move kinds to fire.

\subsection{Verification by program size}\label{sec:eval-per-cat}

\begin{table}[t]
\centering
\caption{Verification results by program size category.}
\label{tab:per-category}
\begin{tabular}{@{}lrrrrr@{}}
\toprule
\textbf{Size} & \textbf{Count} & \textbf{Verified}
  & \textbf{Mean Coords} & \textbf{Mean Props} & \textbf{Mean Time} \\
\midrule
Tiny    & \compTinyCount & \compTinyVerified & \compTinyMeanCoords & \compTinyMeanProps & \compTinyMeanTime \\
Small   & \compSmallCount & \compSmallVerified & \compSmallMeanCoords & \compSmallMeanProps & \compSmallMeanTime \\
Medium  & \compMediumCount & \compMediumVerified & \compMediumMeanCoords & \compMediumMeanProps & \compMediumMeanTime \\
Large   & \compLargeCount & \compLargeVerified & \compLargeMeanCoords & \compLargeMeanProps & \compLargeMeanTime \\
\bottomrule
\end{tabular}
\end{table}

\Cref{tab:per-category} shows that all size categories achieve
complete verification.  Mean coordinates and properties scale
with program size, from \compTinyMeanCoords{} coordinates for
tiny programs to \compLargeMeanCoords{} for large programs.
Verification time remains stable across sizes
(\compTinyMeanTime{} for tiny vs.\ \compLargeMeanTime{} for large),
indicating that the geometric apparatus scales well.

\subsection{Site structure analysis}\label{sec:eval-morph-gt-coord}

\begin{table}[t]
\centering
\caption{Site structure across \compTotalPrograms{} programs.}
\label{tab:morph-gt-coord}
\begin{tabular}{@{}lr@{}}
\toprule
\textbf{Metric} & \textbf{Value} \\
\midrule
Total morphisms           & \compMorphTotal \\
Mean morphisms per site   & \compSiteMeanMorphisms \\
Mean coordinates per site & \compSiteMeanCoords \\
Total functions           & \compSiteTotalFunctions \\
Total classes             & \compSiteTotalClasses \\
Sites constructed         & \compSiteTotalBuilt \\
Mean site build time      & \compSiteMeanBuildMs\,ms \\
\bottomrule
\end{tabular}
\end{table}

\Cref{tab:morph-gt-coord} confirms that morphisms
(\compSiteMeanMorphisms{} per site) exceed coordinates
(\compSiteMeanCoords{} per site) on average, validating the
judgment geometry thesis: for structured programs, the morphism
structure is richer than the coordinate structure.

\subsection{Coverage matrix}\label{sec:eval-coverage}

\begin{table}[t]
\centering
\caption{$\RuleKind \times \TransKind$ coverage matrix.
\checkmark{} = at least one move sequence covers the cell.}
\label{tab:coverage-matrix}
\begin{tabular}{l|ccccc}
\toprule
 & \textbf{FWD} & \textbf{BWD} & \textbf{LAT}
   & \textbf{FIX} & \textbf{INT} \\
\midrule
\textbf{STRUCT} & \checkmark & \checkmark & \checkmark
  & \checkmark & \checkmark \\
\textbf{SEMAN}  & \checkmark & \checkmark & \checkmark
  & \checkmark & \checkmark \\
\textbf{AXIOM}  & \checkmark & \checkmark & \checkmark
  & \checkmark & \checkmark \\
\textbf{DERIV}  & \checkmark & \checkmark & \checkmark
  & \checkmark & \checkmark \\
\bottomrule
\end{tabular}
\end{table}

All 20 cells are covered (\Cref{tab:coverage-matrix}).

\subsection{Comparison with \LEAN{} tactics}\label{sec:eval-lean}

\begin{table}[t]
\centering
\caption{Comparison with \LEAN{} 4 tactics.}
\label{tab:lean-comparison}
\begin{tabular}{@{}lrrr@{}}
\toprule
\textbf{Metric} & \textbf{\LEAN{} 4}
  & \textbf{Sem.\ Moves} & \textbf{Diff} \\
\midrule
Simple lemma        & 3   & 7--12 & +4--9 \\
Medium lemma        & 8   & 12--27 & variable \\
Complex lemma       & 18  & 27--36 & +9--18 \\
Move kinds          & ${\sim}6$ & 13 & +7 \\
Trust manipulation  & No  & Yes & --- \\
Obstruction tracking & No  & Yes & --- \\
Sheaf descent       & No  & Yes & --- \\
\bottomrule
\end{tabular}
\end{table}

\subsection{Timing}\label{sec:eval-timing}

Total elapsed time for \compTotalPrograms{} programs: \compTimeTotal.
Mean verification time per program: \compTimeMean{}
(median \compTimeMedian, std.\ dev.\ \compTimeStdev).
Semantic site construction: \compSiteMeanBuildMs\,ms mean
(\compSiteMaxBuildMs\,ms max).
These sub-millisecond construction times confirm negligible
overhead from the geometric apparatus.

\subsection{Move applicability}\label{sec:eval-failed}

Across \compTotalPrograms{} programs, classical moves always
succeed when their preconditions hold.  Geometric moves may
fail when required structure is absent:
\begin{itemize}[nosep]
  \item \texttt{transport\_along}: requires an admissible
        morphism with a defined pushforward.
  \item \texttt{repair}: requires a refinable covering family.
  \item \texttt{descend}: requires a compatible family over
        a cover.
  \item \texttt{merge}: requires compatible overlap data.
  \item \texttt{restrict\_to}: requires the target coordinate
        to be in the site.
\end{itemize}
In our benchmark, all \compTotalVerified{} programs are
successfully verified, meaning that sufficient geometric
structure exists to complete descent in every case.

\subsection{Descent analysis}\label{sec:eval-descent}

We analyze descent success rates across the four descent strategies.

\begin{table}[t]
\centering
\caption{Descent strategy performance.}
\label{tab:descent-analysis}
\begin{tabular}{@{}lrrrr@{}}
\toprule
\textbf{Strategy} & \textbf{Runs} & \textbf{Successes}
  & \textbf{Rate} & \textbf{Mean (ms)} \\
\midrule
Eager       & \compDescentEagerRuns & \compDescentEagerSuccesses
  & \compDescentEagerRate & \compDescentEagerMeanMs \\
Exhaustive  & \compDescentExhaustiveRuns & \compDescentExhaustiveSuccesses
  & \compDescentExhaustiveRate & \compDescentExhaustiveMeanMs \\
Iterative   & \compDescentIterativeRuns & \compDescentIterativeSuccesses
  & \compDescentIterativeRate & \compDescentIterativeMeanMs \\
Optimistic  & \compDescentOptimisticRuns & \compDescentOptimisticSuccesses
  & \compDescentOptimisticRate & \compDescentOptimisticMeanMs \\
\bottomrule
\end{tabular}
\end{table}

\Cref{tab:descent-analysis} shows that all four descent strategies
achieve \compDescentEagerRate{} success across \compDescentEagerRuns{}
runs.  The strategies agree on all benchmark programs because the
test programs have well-formed covering families where descent
always produces a global section.

In theory, strategies diverge on programs with complex overlap
structures where the order of exploration affects whether a
compatible family is found within the computational budget.
The uniform success rate in our benchmark reflects the structural
regularity of the test programs rather than a general property
of the descent mechanism.

When descent does fail in general (not observed in this benchmark),
the obstruction resides in $\Hcoh{1}(\Site, \J)$, and the
\REPAIR{} move can address such failures by refining the covering
family to separate conflicting regions.

\subsection{Threats to validity}\label{sec:eval-threats}

\paragraph{Internal validity.}
Our evaluation uses synthetic benchmarks generated from
representative Python patterns.  While the \compTotalPrograms{} programs span
four size categories and include both simple and complex structures,
they may not capture all patterns found in production codebases.

\paragraph{External validity.}
All programs are in Python; the morphism distribution and
site structure may differ for languages with
different module systems (e.g., Haskell's type classes, Rust's
trait system).  The benchmark programs range up to
\compLargeMeanLines{} mean lines for the large category;
larger programs may exhibit different morphism ratios.

\paragraph{Construct validity.}
We report morphism distribution, site structure, and descent
success rates rather than subjective metrics.
An alternative weighting by trust gain or evidence quality might
yield different conclusions about the importance of geometric
apparatus.

\paragraph{Conclusion validity.}
Our sample of \compTotalPrograms{} programs
(\compTinyCount{} tiny, \compSmallCount{} small,
\compMediumCount{} medium, \compLargeCount{} large)
is modest.  The \compOverallAccuracy{} verification rate and
uniform descent success reflect the structural regularity of
the benchmark rather than a general property of all Python programs.


%% =====================================================================
\section{Comparison with Existing Tactic Systems}\label{sec:comparison}
%% =====================================================================

\subsection{Feature matrix}

\begin{table}[t]
\centering
\caption{Feature comparison of tactic frameworks.}
\label{tab:feature-matrix}
\begin{tabular}{l|cccc}
\toprule
\textbf{Feature} & \textbf{\LEAN{} 4} & \textbf{Ltac2}
  & \textbf{Meta-\Fstar} & \textbf{Sem.\ Moves} \\
\midrule
Typed tactics          & \checkmark & \checkmark & \checkmark & \checkmark \\
Sequential comp.       & \checkmark & \checkmark & \checkmark & \checkmark \\
Parallel comp.         &            &            &            & \checkmark \\
Backtracking           & \checkmark & \checkmark &            & \checkmark \\
SMT integration        &            &            & \checkmark & \checkmark \\
Trust manipulation     &            &            &            & \checkmark \\
Coordinate-aware       &            &            &            & \checkmark \\
Cover-aware            &            &            &            & \checkmark \\
Obstruction tracking   &            &            &            & \checkmark \\
Sheaf descent          &            &            &            & \checkmark \\
Repair generation      &            &            &            & \checkmark \\
\bottomrule
\end{tabular}
\end{table}

\subsection{Tactic simulation analysis}

Of the 13 move kinds, the 6 classical moves can be simulated:
\texttt{compose} $\leftrightarrow$ \texttt{apply/exact},
\texttt{split} $\leftrightarrow$ \texttt{split/cases},
\texttt{strengthen} $\leftrightarrow$ \texttt{assumption},
\texttt{weaken} $\leftrightarrow$ \texttt{sorry/admit},
\texttt{add\_obligation} $\leftrightarrow$ \texttt{have/suffices},
\texttt{discharge\_obligation} $\leftrightarrow$ \texttt{exact}.

The 7 geometric moves cannot be simulated
(\Cref{thm:inexpressible}).  We now give a detailed proof
sketch for each geometric move.

\subsubsection{Inexpressibility of \texttt{restrict\_to}}

The \texttt{restrict\_to} move narrows the proof state to a
sub-coordinate, restricting all components (evidence, obligations,
obstructions, trust) simultaneously.  In \LEAN, one can focus on
a subgoal via \texttt{case} or \texttt{show}, but these operate
on the proof \emph{term}, not on a coordinate in a semantic site.
There is no \LEAN{} tactic that restricts trust annotations,
obstruction presheaves, or covering families to a sub-site.
The state components that \texttt{restrict\_to} modifies
(site topology, trust map, obstruction presheaf) have no
representation in the \LEAN{} metavariable context.

\subsubsection{Inexpressibility of \texttt{merge}}

The \texttt{merge} move computes a pushout of two local proof
states along their overlap.  This requires:
(i) an overlap intersection $c_1 \cap c_2$ in the site;
(ii) a compatibility check on evidence over the overlap;
(iii) a conservative join of trust levels.
\LEAN's \texttt{simp} and \texttt{ring} can combine lemmas,
but they operate within a single proof context, not across
two independent proof states with their own trust levels
and obstruction data.  The merge operation is a categorical
colimit, and \LEAN's tactic monad does not support colimit
computation over proof contexts.

\subsubsection{Inexpressibility of \texttt{transport\_along}}

Transport applies a pushforward functor $f_*$ along a site
morphism $f: c_1 \to c_2$.  While \LEAN's \texttt{rw}
tactic rewrites along an equality $a = b$, it does not
transport along a categorical morphism between different
coordinates.  The pushforward requires tracking how evidence,
obligations, and trust transform covariantly along $f$---a
functorial operation that has no analogue in term rewriting.
Moreover, \texttt{rw} does not modify trust levels, which
$f_*$ does by the monotonicity property.

\subsubsection{Inexpressibility of \texttt{add\_obstruction}
and \texttt{resolve\_obstruction}}

These moves modify the obstruction presheaf $\obstr$, a
component that does not exist in any existing tactic framework.
Adding an obstruction records a cohomological obstruction at
a coordinate; resolving one removes it by providing evidence.
\LEAN's closest analogue is \texttt{sorry}, which admits a
goal without evidence, but \texttt{sorry} does not record
an obstruction---it silently assumes the goal is true.
The obstruction presheaf enables the system to track
\emph{why} a proof attempt failed and \emph{where} the
failure occurred in the site topology, information that
\texttt{sorry} discards.

\subsubsection{Inexpressibility of \texttt{repair}}

Repair modifies the Grothendieck topology itself---refining
covering families to eliminate obstructions.  No tactic in
\LEAN, Ltac2, or Meta-\Fstar{} can modify the topology of
the proof state because the proof state has no topology.
The closest analogue in \LEAN{} is changing the set of
hypotheses via \texttt{have} or \texttt{obtain}, but these
add hypotheses to a single proof context rather than refining
a covering family over a topological space.

\subsubsection{Inexpressibility of \texttt{descend}}

Descent computes a global section from compatible local data
by the sheaf axiom.  This requires: (i) a covering family;
(ii) local sections satisfying the cocycle condition on
overlaps; (iii) the unique gluing provided by the sheaf
property.  \LEAN's proof state has no covering families, no
overlap computations, and no sheaf axiom.  While one could
manually encode the descent computation as a \LEAN{} proof,
this would require the user to implement sheaf theory inside
\LEAN---the descent operation would not be a \emph{tactic}
but rather a \emph{theorem}, and the proof search automation
provided by the \texttt{descend} move would be lost.

\subsection{Expressiveness ordering}

\[
  \LEAN \;\equiv_{\mathrm{exp}}\; \text{Ltac2}
  \;\equiv_{\mathrm{exp}}\; \text{Meta-\Fstar}
  \;\subsetneq_{\mathrm{exp}}\; \text{Semantic Moves}
\]
The gap is the 7 geometric move kinds.


%% =====================================================================
\section{Related Work and Conclusion}\label{sec:conclusion}
%% =====================================================================

\subsection{Related work}

\paragraph{Tactic frameworks.}
\LEAN{} 4~\cite{deMoura2021Lean4}, Ltac2~\cite{Pedrot2019Ltac2},
and Meta-\Fstar{}~\cite{MartinezEtAl2019MetaFStar} are closest.
Mtac2~\cite{Kaiser2018Mtac2} and Eisbach~\cite{Matichuk2016Eisbach}
add typed meta-programming and method languages but remain term-based.

\paragraph{Proof search and automation.}
Hammer tools~\cite{Blanchette2016Hammers} erase geometric structure.
ML-guided search~\cite{Polu2020GPTf,Yang2024LeanDojo} could benefit
from 7 additional geometric move kinds as an action space.

\paragraph{Sheaf-theoretic methods.}
Abramsky et~al.~\cite{Abramsky2011Sheaves} use sheaves for
contextuality in quantum mechanics.  Our application to proof search
is the first use of Grothendieck topologies for tactic design.

\paragraph{Compositional verification.}
Hoare logic, rely-guarantee, and separation logic compose proofs
within a fixed calculus.  Semantic moves add topological composition
via covering families.

\paragraph{Proof repair.}
Ringer et~al.~\cite{Ringer2021ProofRepair} repair proof terms.
Our \texttt{repair} move modifies the site structure, a complementary
approach.

\paragraph{Program synthesis.}
Program synthesis systems such as Sketch~\cite{Solar-Lezama2008Sketch}
and Rosette~\cite{Torlak2014Rosette} construct programs from
specifications using counterexample-guided inductive synthesis (CEGIS).
While CEGIS iteratively refines candidates, it operates on a flat
search space of program terms without geometric structure.
Semantic moves could augment synthesis by providing a richer
action space: the \CONSTRUCT{} move generates candidate
local sections, while \GLUE{} and \ESCALATE{} compose them
across module boundaries---operations that CEGIS cannot express
because it lacks a notion of covering families or overlap
compatibility.

\paragraph{Proof planning.}
Bundy's proof planning framework~\cite{Bundy1988ProofPlanning}
organizes proof search around high-level \emph{methods} that
encode common proof patterns.  Methods have preconditions,
postconditions, and effects---similar to our semantic moves.
However, proof planning methods operate on proof terms and
do not manipulate trust levels, site structure, or
obstructions.  Our move taxonomy can be viewed as an
extension of proof planning to a geometric setting, where
methods additionally modify the topology of the proof state
and where the control layer uses obstruction norms rather
than heuristic method selection.

\paragraph{Rippling.}
Rippling~\cite{Bundy2005Rippling} is a heuristic for
guiding rewriting in inductive proofs by annotating terms
with \emph{wave fronts} and directing rewrites toward
the induction hypothesis.  While rippling provides directed
search within a single proof context, semantic moves provide
directed search across \emph{multiple} proof contexts linked
by site morphisms.  The \TRANSPORT{} move is the geometric
analogue of rippling: both direct proof progress along a
structural relationship, but transport operates along
categorical morphisms rather than syntactic wave fronts.

\paragraph{Strategy languages.}
Strategy languages such as Stratego~\cite{Visser2004Stratego}
and Elan~\cite{Borovansky1998Elan} provide programmable
term-rewriting strategies with combinators for sequencing,
choice, and recursion.  These languages are Turing-complete
over term transformations but cannot express geometric
operations: there is no combinator for restricting to a
sub-site, gluing along an overlap, or performing sheaf
descent.  Semantic moves extend the strategy language paradigm
from term rewriting to site-level proof manipulation.

\paragraph{Proof by analogy.}
Proof by analogy~\cite{Melis1999ProofAnalogy} transfers proof
structure from a solved problem to a new one via a mapping
between their representations.  Our \TRANSPORT{} move
formalizes a restricted form of proof analogy: transporting
evidence along a site morphism $f: c_1 \to c_2$ is analogous
to mapping proof structure from $c_1$ to $c_2$.  However,
\TRANSPORT{} is functorial and trust-aware, properties that
the informal analogy framework of Melis lacks.  Moreover,
semantic moves compose analogies via the move algebra
(\Cref{sec:move-algebra}), enabling iterated transport
along chains of morphisms.

\subsection{Conclusion}

We introduced semantic moves---typed, composable proof-search
operations acting on the full geometric proof state.  Of 13
move kinds, 7 are geometric and have no \LEAN/\Coq/\Fstar{}
analogue.

Evaluation on \compTotalPrograms{} programs across four size categories shows:
\begin{itemize}[nosep]
  \item \compMorphTotal{} morphisms classified as restriction
        (\compMorphRestrictionPct), transport (\compMorphTransportPct),
        and inclusion (\compMorphInclusionPct).
  \item Morphisms per site (\compSiteMeanMorphisms) exceed
        coordinates per site (\compSiteMeanCoords) on average.
  \item All 20 rule--transition cells are covered.
  \item All four descent strategies achieve \compDescentEagerRate{} success.
  \item The 10 named moves (\VERIFY, \CONSTRUCT, \NORMALIZE,
        \GLUE, \RESTRICT, \TRANSPORT, \REPAIR, \PROMOTE,
        \CHALLENGE, \ESCALATE) provide a complete orchestration
        vocabulary.
\end{itemize}

Future work: larger programs where \texttt{transport\_along}
and \texttt{repair} succeed more frequently, ML-guided
move selection, and a \LEAN{} bridge for classical-move export.


%% =====================================================================
\appendix
\section{Representative Programs}\label{app:programs}
%% =====================================================================

\Cref{tab:representative} summarises the real evaluation data
across our four size categories, with coordinate and morphism
counts taken directly from the comprehensive benchmark.

\begin{table}[htbp]
\centering
\caption{Verification results by program-size category (real data).}
\label{tab:representative}
\small
\begin{tabular}{@{}lrrrrl@{}}
\toprule
\textbf{Category} & \textbf{$n$} & \textbf{Mean~C} & \textbf{Mean~P}
  & \textbf{Mean~$t$} & \textbf{Verified} \\
\midrule
Tiny ($\sim$5\,LOC)
  & \compTinyCount & \compTinyMeanCoords & \compTinyMeanProps
  & \compTinyMeanTime & \compTinyVerified/\compTinyCount \\
Small ($\sim$16\,LOC)
  & \compSmallCount & \compSmallMeanCoords & \compSmallMeanProps
  & \compSmallMeanTime & \compSmallVerified/\compSmallCount \\
Medium ($\sim$45\,LOC)
  & \compMediumCount & \compMediumMeanCoords & \compMediumMeanProps
  & \compMediumMeanTime & \compMediumVerified/\compMediumCount \\
Large ($\sim$79\,LOC)
  & \compLargeCount & \compLargeMeanCoords & \compLargeMeanProps
  & \compLargeMeanTime & \compLargeVerified/\compLargeCount \\
\midrule
\textbf{Total}
  & \compTotalPrograms & \compCoordsMean & \compPropsMean
  & \compTimeMean & \compTotalVerified/\compTotalPrograms \\
\bottomrule
\end{tabular}
\par\medskip
\footnotesize
C = coordinate count, P = proposition count, $t$ = wall-clock time.
All data from \texttt{comprehensive-data.tex} (auto-generated).
\end{table}

Several patterns are visible.  The morphism type distribution across
the full benchmark is \compMorphRestrictionPct{} restriction,
\compMorphTransportPct{} transport, and \compMorphInclusionPct{}
inclusion morphisms (out of \compMorphTotal{} total).  Larger
programs produce richer site structure (up to \compCoordsMax{}
coordinates and \compPropsMax{} propositions for 79-line programs)
while maintaining full verification.  Timing is dominated by the SMT
solver and remains roughly constant across size categories
($\sigma = $ \compTimeStdev).
Modules (\texttt{protocol\_handler}, \texttt{test\_framework})
have the most morphisms and the highest geometric fractions,
driven by cross-module transport and descent operations.


%% =====================================================================
%% Bibliography
%% =====================================================================

\begin{thebibliography}{20}

\bibitem{deMoura2021Lean4}
L.~de~Moura et~al.
\newblock The {Lean} 4 theorem prover.
\newblock In \emph{CADE-28}, Springer LNCS, 2021.

\bibitem{Avigad2023Hitchhiker}
J.~Avigad et~al.
\newblock The hitchhiker's guide to logical verification, 2023.

\bibitem{Pedrot2019Ltac2}
P.-M.~P\'{e}drot.
\newblock Ltac2: Tactical warfare.
\newblock In \emph{CoqPL Workshop}, 2019.

\bibitem{MartinezEtAl2019MetaFStar}
G.~Mart\'{i}nez et~al.
\newblock Meta-{F}$^{\star}$: Proof automation with SMT, tactics,
and metaprograms.
\newblock In \emph{ESOP}, Springer LNCS, 2019.

\bibitem{JuGeo1Sites}
H.~Young.
\newblock Semantic sites for program verification.
\newblock Judgment Geometry Paper~1, 2024.

\bibitem{JuGeo2Algebra}
H.~Young.
\newblock The judgment algebra.
\newblock Judgment Geometry Paper~2, 2024.

\bibitem{JuGeo3Descent}
H.~Young.
\newblock Descent and obstructions.
\newblock Judgment Geometry Paper~3, 2024.

\bibitem{JuGeo4Trust}
H.~Young.
\newblock Trust algebras for multi-source verification.
\newblock Judgment Geometry Paper~4, 2024.

\bibitem{Kaiser2018Mtac2}
J.-O.~Kaiser et~al.
\newblock Mtac2: Typed tactics for backward reasoning in {Coq}.
\newblock In \emph{ICFP}, ACM, 2018.

\bibitem{Matichuk2016Eisbach}
D.~Matichuk, T.~Murray, and M.~Wenzel.
\newblock Eisbach: A proof method language for {Isabelle}.
\newblock \emph{J.\ Autom.\ Reasoning}, 56(3):261--282, 2016.

\bibitem{Blanchette2016Hammers}
J.~C.~Blanchette et~al.
\newblock Hammering towards {QED}.
\newblock \emph{J.\ Formalized Reasoning}, 9(1):101--148, 2016.

\bibitem{Polu2020GPTf}
S.~Polu and I.~Sutskever.
\newblock Generative language modeling for automated theorem proving.
\newblock \emph{arXiv:2009.03393}, 2020.

\bibitem{Yang2024LeanDojo}
K.~Yang et~al.
\newblock {LeanDojo}: Theorem proving with retrieval-augmented LMs.
\newblock In \emph{NeurIPS}, 2023.

\bibitem{MacLaneMoerdijk1994}
S.~Mac~Lane and I.~Moerdijk.
\newblock \emph{Sheaves in Geometry and Logic}.
\newblock Springer, 1994.

\bibitem{SGA4}
M.~Artin, A.~Grothendieck, and J.-L.~Verdier.
\newblock \emph{SGA~4}.
\newblock Springer LNM 269/270/305, 1972--1973.

\bibitem{Abramsky2011Sheaves}
S.~Abramsky and A.~Brandenburger.
\newblock The sheaf-theoretic structure of non-locality.
\newblock \emph{New J.\ Phys.}, 13:113036, 2011.

\bibitem{Ringer2021ProofRepair}
T.~Ringer et~al.
\newblock Proof repair across type equivalences.
\newblock In \emph{PLDI}, ACM, 2021.

\bibitem{Solar-Lezama2008Sketch}
A.~Solar-Lezama.
\newblock Program synthesis by sketching.
\newblock PhD thesis, UC Berkeley, 2008.

\bibitem{Torlak2014Rosette}
E.~Torlak and R.~Bod\'{i}k.
\newblock A lightweight symbolic virtual machine for solver-aided
host languages.
\newblock In \emph{PLDI}, ACM, 2014.

\bibitem{Bundy1988ProofPlanning}
A.~Bundy.
\newblock The use of explicit plans to guide inductive proofs.
\newblock In \emph{CADE}, Springer LNCS, 1988.

\bibitem{Bundy2005Rippling}
A.~Bundy.
\newblock Rippling: Meta-level guidance for mathematical reasoning.
\newblock Cambridge University Press, 2005.

\bibitem{Visser2004Stratego}
E.~Visser.
\newblock Program transformation with {Stratego/XT}.
\newblock In \emph{Domain-Specific Program Generation}, Springer LNCS, 2004.

\bibitem{Borovansky1998Elan}
P.~Borovansk\'{y} et~al.
\newblock {ELAN} from a rewriting logic point of view.
\newblock \emph{Theoretical Computer Science}, 285(2):155--185, 2002.

\bibitem{Melis1999ProofAnalogy}
E.~Melis.
\newblock The heuristic solver: Analogy in proof search.
\newblock In \emph{Advances in Analogy Research}, New Bulgarian University, 1998.

\end{thebibliography}

\end{document}
